\documentclass[../main.tex]{subfiles}
\begin{document}

\chapter{分散ファイルシステム}
\lessoninfo{
  video={P4L2},
  textbook={OSTEP \cite{ostep} ch.~49--50；Silberschatz ら \cite{silberschatz2018} ch.~15, 19；Nelson ら \cite{nelson1988}},
  keywords={分散ファイルシステム，レプリケーション，パーティショニング，アップロード・ダウンロードモデル，ステートレスサーバ，ステートフルサーバ，UNIXセマンティクス，セッションセマンティクス，リース，NFS，ファイルハンドル，Sprite，遅延書き込み，並行書き込み共有},
}

第11章の仮想ファイルシステムは，アプリケーションが使うファイルがどこに
あるかをアプリケーションから隠す．第13章の遠隔手続き呼び出しは，プロセスが
別のマシンで操作を実行させることを可能にする．分散ファイルシステムはこの
2つを組み合わせたものであり，ネットワークで結ばれたマシン上にファイルを
格納し，通常のファイルシステムインタフェースを使うクライアントにそれらを
提供する．この章では，そうしたシステムの設計空間---クライアントとサーバの
間で作業をどう分けるか，サーバがどのような状態を保持するか，キャッシュ
された複製をどう一貫させるか，ファイルを共有したときクライアントが何を
観測するか---を調べ，続いて2つの具体的な設計，すなわち NFS と，ファイルが
実際にどう使われるかの測定から設計を導いた Sprite を扱う．

\section{分散ファイルシステム}
\label{sec:l14-dfs}

オペレーティングシステムは，記憶装置の詳細をファイルシステムインタフェース
の背後に隠す．\source{P4L2，レッスンの概要，視覚的な比喩，分散ファイル
システム，DFS のモデル；Silberschatz ら ch.~19．}その下では，仮想ファイル
システムが各操作を，ファイルが存在するファイルシステムへ渡す．そのファイル
システムは，ファイルをローカルなデバイスに置くこともあれば，ネットワーク上の
他のマシンに置くこともある．

\begin{definition}[分散ファイルシステム]
  \term[ぶんさんふぁいるしすてむ]{分散ファイルシステム}[distributed file system]%
  \index{DFS|see{分散ファイルシステム}}（DFS）とは，ネットワークで結ばれた
  マシン上に格納され，そこから提供されるファイルからなるファイルシステム
  であって，プロセスは，ファイルが自分のマシンにあるかどうかにかかわらず，
  通常のファイルシステムインタフェースを通じてアクセスする．
\end{definition}

DFS を特徴づける性質は3つある．
\begin{itemize}
  \item 明確に定義されたインタフェースを通じてアクセスされる．アプリケー
    ションはローカルなファイルに対するのと同じシステムコールを使い，VFS は
    リモートファイルに対する各操作を DFS へ渡す．
  \item 一貫した状態を保たなければならない．複数のクライアントが同じ
    ファイルを使うとき，DFS はどのクライアントがどのファイルにアクセス
    しているかを追跡し，更新を伝播し，キャッシュされた複製の整合を保た
    なければならない．
  \item 複数の分散モデルを組み合わせうる．ファイルは複数のサーバに複製する
    ことも分割することもでき，クライアントとサーバの役割が別々のマシンに
    存在する必要もない．
\end{itemize}

分散のモデルを \cref{fig:l14-models} に示す．
\begin{description}
  \item[クライアントとサーバが別のマシンにある場合] 単一のファイルサーバが
    すべてのファイルを格納し，他のマシン上のクライアントがネットワーク越し
    にそれらにアクセスする．
  \item[分散ファイルサーバ] ファイルサービス自体が複数のマシンで動作する．
    \emph{複製された}ファイルサーバでは，すべてのマシンがすべてのファイルを
    保持し，\emph{分割された}ファイルサーバでは，各マシンがファイルの一部を
    保持する．両者は \cref{sec:l14-replication} で比較する．
  \item[ピア] すべてのマシンがファイルを格納し，かつ提供する．
\end{description}

最後のモデルでは，クライアントとサーバの区別は曖昧になる．このような構成を
\term[ぴあつーぴあ]{ピアツーピア}[peer-to-peer]と呼ぶ．どのマシンも，
ファイルサービスを利用すると同時に提供するピアである．各ピアはファイルの
一部を格納し，要求---典型的には自分が格納しているファイルに対する要求---を
処理することで負荷の一部を担う．

\begin{figure}[htbp]
  \centering
  % DFS models: (a) client and server on different machines, (b) replicated
% file servers, (c) partitioned file servers, (d) peers.
% Usage: % DFS models: (a) client and server on different machines, (b) replicated
% file servers, (c) partitioned file servers, (d) peers.
% Usage: % DFS models: (a) client and server on different machines, (b) replicated
% file servers, (c) partitioned file servers, (d) peers.
% Usage: \input{figures/l14-models}   (width ~118 mm)
\begin{tikzpicture}[x=1mm, y=1mm, font=\scriptsize\sffamily,
  cl/.style={draw, fill=white, minimum width=12mm, minimum height=5mm, inner sep=1pt},
  sv/.style={draw, fill=ln-box, minimum width=15mm, minimum height=8mm, inner sep=1pt, align=center},
  pe/.style={draw, fill=ln-accent!12, minimum width=15mm, minimum height=8mm, inner sep=1pt, align=center},
  net/.style={line width=1.2pt, ln-muted},
  lab/.style={font=\scriptsize\sffamily, text=ln-muted},
  ttl/.style={font=\footnotesize\sffamily\bfseries, anchor=west}]
  % (a) one server
  \begin{scope}[shift={(0,33)}]
    \node[ttl] at (0,27) {\iflangja{(a) 単一のファイルサーバ}{(a) single file server}};
    \draw[net] (3,13) -- (53,13);
    \foreach \x in {10,28,46} {
      \node[cl] at (\x,20) {\iflangja{クライアント}{client}};
      \draw (\x,17.5) -- (\x,13);
    }
    \node[lab, anchor=south] at (19,13.3) {\iflangja{ネットワーク}{network}};
    \node[sv] at (28,4) {\iflangja{サーバ}{server}\\A--F};
    \draw (28,8) -- (28,13);
  \end{scope}
  % (b) replicated
  \begin{scope}[shift={(62,33)}]
    \node[ttl] at (0,27) {\iflangja{(b) 複製されたファイルサーバ}{(b) replicated file servers}};
    \draw[net] (3,13) -- (53,13);
    \foreach \x in {10,28,46} {
      \node[cl] at (\x,20) {\iflangja{クライアント}{client}};
      \draw (\x,17.5) -- (\x,13);
      \node[sv] at (\x,4) {\iflangja{サーバ}{server}\\A--F};
      \draw (\x,8) -- (\x,13);
    }
  \end{scope}
  % (c) partitioned
  \begin{scope}[shift={(0,0)}]
    \node[ttl] at (0,27) {\iflangja{(c) 分割されたファイルサーバ}{(c) partitioned file servers}};
    \draw[net] (3,13) -- (53,13);
    \foreach \x/\f in {10/{A, B},28/{C, D},46/{E, F}} {
      \node[cl] at (\x,20) {\iflangja{クライアント}{client}};
      \draw (\x,17.5) -- (\x,13);
      \node[sv] at (\x,4) {\iflangja{サーバ}{server}\\\f};
      \draw (\x,8) -- (\x,13);
    }
  \end{scope}
  % (d) peers
  \begin{scope}[shift={(62,0)}]
    \node[ttl] at (0,27) {\iflangja{(d) ピア}{(d) peers}};
    \draw[net] (3,19) -- (53,19);
    \foreach \x/\f in {10/{A, B},28/{C, D},46/{E, F}} {
      \node[pe] at (\x,10) {\iflangja{ピア}{peer}\\\f};
      \draw (\x,14) -- (\x,19);
    }
    \node[lab, anchor=north] at (28,5)
      {\iflangja{各ピアがクライアントでありサーバでもある}{each peer is both client and server}};
  \end{scope}
\end{tikzpicture}
   (width ~118 mm)
\begin{tikzpicture}[x=1mm, y=1mm, font=\scriptsize\sffamily,
  cl/.style={draw, fill=white, minimum width=12mm, minimum height=5mm, inner sep=1pt},
  sv/.style={draw, fill=ln-box, minimum width=15mm, minimum height=8mm, inner sep=1pt, align=center},
  pe/.style={draw, fill=ln-accent!12, minimum width=15mm, minimum height=8mm, inner sep=1pt, align=center},
  net/.style={line width=1.2pt, ln-muted},
  lab/.style={font=\scriptsize\sffamily, text=ln-muted},
  ttl/.style={font=\footnotesize\sffamily\bfseries, anchor=west}]
  % (a) one server
  \begin{scope}[shift={(0,33)}]
    \node[ttl] at (0,27) {\iflangja{(a) 単一のファイルサーバ}{(a) single file server}};
    \draw[net] (3,13) -- (53,13);
    \foreach \x in {10,28,46} {
      \node[cl] at (\x,20) {\iflangja{クライアント}{client}};
      \draw (\x,17.5) -- (\x,13);
    }
    \node[lab, anchor=south] at (19,13.3) {\iflangja{ネットワーク}{network}};
    \node[sv] at (28,4) {\iflangja{サーバ}{server}\\A--F};
    \draw (28,8) -- (28,13);
  \end{scope}
  % (b) replicated
  \begin{scope}[shift={(62,33)}]
    \node[ttl] at (0,27) {\iflangja{(b) 複製されたファイルサーバ}{(b) replicated file servers}};
    \draw[net] (3,13) -- (53,13);
    \foreach \x in {10,28,46} {
      \node[cl] at (\x,20) {\iflangja{クライアント}{client}};
      \draw (\x,17.5) -- (\x,13);
      \node[sv] at (\x,4) {\iflangja{サーバ}{server}\\A--F};
      \draw (\x,8) -- (\x,13);
    }
  \end{scope}
  % (c) partitioned
  \begin{scope}[shift={(0,0)}]
    \node[ttl] at (0,27) {\iflangja{(c) 分割されたファイルサーバ}{(c) partitioned file servers}};
    \draw[net] (3,13) -- (53,13);
    \foreach \x/\f in {10/{A, B},28/{C, D},46/{E, F}} {
      \node[cl] at (\x,20) {\iflangja{クライアント}{client}};
      \draw (\x,17.5) -- (\x,13);
      \node[sv] at (\x,4) {\iflangja{サーバ}{server}\\\f};
      \draw (\x,8) -- (\x,13);
    }
  \end{scope}
  % (d) peers
  \begin{scope}[shift={(62,0)}]
    \node[ttl] at (0,27) {\iflangja{(d) ピア}{(d) peers}};
    \draw[net] (3,19) -- (53,19);
    \foreach \x/\f in {10/{A, B},28/{C, D},46/{E, F}} {
      \node[pe] at (\x,10) {\iflangja{ピア}{peer}\\\f};
      \draw (\x,14) -- (\x,19);
    }
    \node[lab, anchor=north] at (28,5)
      {\iflangja{各ピアがクライアントでありサーバでもある}{each peer is both client and server}};
  \end{scope}
\end{tikzpicture}
   (width ~118 mm)
\begin{tikzpicture}[x=1mm, y=1mm, font=\scriptsize\sffamily,
  cl/.style={draw, fill=white, minimum width=12mm, minimum height=5mm, inner sep=1pt},
  sv/.style={draw, fill=ln-box, minimum width=15mm, minimum height=8mm, inner sep=1pt, align=center},
  pe/.style={draw, fill=ln-accent!12, minimum width=15mm, minimum height=8mm, inner sep=1pt, align=center},
  net/.style={line width=1.2pt, ln-muted},
  lab/.style={font=\scriptsize\sffamily, text=ln-muted},
  ttl/.style={font=\footnotesize\sffamily\bfseries, anchor=west}]
  % (a) one server
  \begin{scope}[shift={(0,33)}]
    \node[ttl] at (0,27) {\iflangja{(a) 単一のファイルサーバ}{(a) single file server}};
    \draw[net] (3,13) -- (53,13);
    \foreach \x in {10,28,46} {
      \node[cl] at (\x,20) {\iflangja{クライアント}{client}};
      \draw (\x,17.5) -- (\x,13);
    }
    \node[lab, anchor=south] at (19,13.3) {\iflangja{ネットワーク}{network}};
    \node[sv] at (28,4) {\iflangja{サーバ}{server}\\A--F};
    \draw (28,8) -- (28,13);
  \end{scope}
  % (b) replicated
  \begin{scope}[shift={(62,33)}]
    \node[ttl] at (0,27) {\iflangja{(b) 複製されたファイルサーバ}{(b) replicated file servers}};
    \draw[net] (3,13) -- (53,13);
    \foreach \x in {10,28,46} {
      \node[cl] at (\x,20) {\iflangja{クライアント}{client}};
      \draw (\x,17.5) -- (\x,13);
      \node[sv] at (\x,4) {\iflangja{サーバ}{server}\\A--F};
      \draw (\x,8) -- (\x,13);
    }
  \end{scope}
  % (c) partitioned
  \begin{scope}[shift={(0,0)}]
    \node[ttl] at (0,27) {\iflangja{(c) 分割されたファイルサーバ}{(c) partitioned file servers}};
    \draw[net] (3,13) -- (53,13);
    \foreach \x/\f in {10/{A, B},28/{C, D},46/{E, F}} {
      \node[cl] at (\x,20) {\iflangja{クライアント}{client}};
      \draw (\x,17.5) -- (\x,13);
      \node[sv] at (\x,4) {\iflangja{サーバ}{server}\\\f};
      \draw (\x,8) -- (\x,13);
    }
  \end{scope}
  % (d) peers
  \begin{scope}[shift={(62,0)}]
    \node[ttl] at (0,27) {\iflangja{(d) ピア}{(d) peers}};
    \draw[net] (3,19) -- (53,19);
    \foreach \x/\f in {10/{A, B},28/{C, D},46/{E, F}} {
      \node[pe] at (\x,10) {\iflangja{ピア}{peer}\\\f};
      \draw (\x,14) -- (\x,19);
    }
    \node[lab, anchor=north] at (28,5)
      {\iflangja{各ピアがクライアントでありサーバでもある}{each peer is both client and server}};
  \end{scope}
\end{tikzpicture}

  \caption{DFS のモデル．(a) 単一のサーバがすべてのファイルを格納する．
    (b) 複製されたファイルサーバでは，どのサーバもすべてのファイルを保持
    する．(c) 分割されたファイルサーバでは，各サーバがファイルの一部を
    保持する．(d) すべてのマシンがピアであり，ファイルを格納して提供する．}
  \label{fig:l14-models}
\end{figure}

\section{遠隔ファイルサービス}
\label{sec:l14-remote}

サーバが1台であっても，リモートファイルにアクセスする作業をクライアントと
サーバの間で分ける方法は大きく異なりうる．\source{P4L2，遠隔ファイル
サービス：両極端，遠隔ファイルサービス：折衷案；OSTEP ch.~49--50；
Silberschatz ら ch.~19．}設計空間を画する両極端が2つある
（\cref{fig:l14-remote}）．

\begin{definition}[アップロード・ダウンロードモデル]
  \term[あっぷろーどだうんろーどもでる]{アップロード・ダウンロードモデル}%
  [upload/download model]では，ファイルにアクセスするクライアントは
  サーバからファイル全体をダウンロードし，すべての操作をローカルな複製に
  対して行い，終わったらファイル全体をサーバへアップロードする．
\end{definition}

FTP によるファイル転送や，Subversion のようなバージョン管理システムは
この方式で動作する．読み出しも書き込みもクライアント側でローカルに行われる
ので高速である．しかし，クライアントがファイルのごく一部しか必要としない
場合でもファイル全体が転送され，さらにサーバは制御を手放す．いったん
ファイルがダウンロードされると，サーバは，クライアントがそれをどう扱うかも，
変更されたファイルがいつアップロードされるかも知りえない．

\begin{definition}[遠隔アクセスモデル]
  \term[えんかくあくせすもでる]{遠隔アクセスモデル}[remote access model]，
  すなわち真の遠隔ファイルアクセスでは，ファイルはサーバ上に留まり，
  ファイルに対するあらゆる操作---すべての読み出しと書き込み---がサーバへ
  送られる．クライアント側でローカルに行われることは何もない．
\end{definition}

すべてのアクセスがサーバを経由するので，ファイルアクセスは集中化され，
一貫性は，複数のプロセスで共有されるローカルなファイルシステムと同じくらい
容易に理解できる．しかし，すべてのファイル操作がネットワークの往復と
サーバでの処理の代価を払い，またサーバがすべてのクライアントのすべての操作を
実行するので，DFS が支えられるクライアント数はサーバの処理能力によって
制限される．

\begin{figure}[htbp]
  \centering
  % Remote file service: (a) upload/download model, (b) remote access model,
% (c) the compromise in which clients cache some blocks.
% Usage: % Remote file service: (a) upload/download model, (b) remote access model,
% (c) the compromise in which clients cache some blocks.
% Usage: % Remote file service: (a) upload/download model, (b) remote access model,
% (c) the compromise in which clients cache some blocks.
% Usage: \input{figures/l14-remote}   (width ~118 mm)
\begin{tikzpicture}[x=1mm, y=1mm, font=\scriptsize\sffamily,
  >={Stealth[length=1.5mm]},
  machine/.style={draw, fill=white, rounded corners=0.8mm},
  lab/.style={font=\scriptsize\sffamily, align=center},
  who/.style={font=\scriptsize\sffamily\itshape, text=ln-muted, anchor=north west},
  ttl/.style={font=\footnotesize\sffamily\bfseries, anchor=west}]
  % \lblocks{x0}{y0}{list of block indices to fill}{fill colour}
  \def\lblocks#1#2#3#4{%
    \foreach \i in {0,...,5} { \draw[fill=white] (#1+\i*4,#2) rectangle ++(4,4); }
    \foreach \i in {#3} { \draw[fill=#4] (#1+\i*4,#2) rectangle ++(4,4); }}
  \foreach \px/\t in {0/{\iflangja{(a) アップロード・ダウンロード}{(a) upload/download}},
                      40/{\iflangja{(b) 遠隔アクセス}{(b) remote access}},
                      80/{\iflangja{(c) 折衷案}{(c) compromise}}} {
    \node[ttl] at (\px,44) {\t};
    \draw[machine] (\px+1,26) rectangle (\px+37,39);
    \node[who] at (\px+1.5,38.6) {\iflangja{クライアント}{client}};
    \draw[machine, fill=ln-box] (\px+1,0) rectangle (\px+37,11);
    \node[who] at (\px+1.5,10.6) {\iflangja{サーバ}{server}};
    \lblocks{\px+7}{1.8}{0,...,5}{ln-muted!45}
  }
  % (a) whole file at the client
  \lblocks{7}{28}{0,...,5}{ln-accent!45}
  \draw[->, thick] (9,11.5) -- node[lab, right, pos=0.7] {\iflangja{ダウンロード}{download}} (9,25.5);
  \draw[->, thick] (29,25.5) -- node[lab, left, pos=0.7] {\iflangja{アップロード}{upload}} (29,11.5);
  % (b) nothing at the client, every operation to the server
  \node[lab, text=ln-muted] at (59,30) {\iflangja{ローカルにはデータなし}{no file data kept locally}};
  \foreach \x in {50,54,58} { \draw[->] (\x,25.5) -- (\x,11.5); }
  \foreach \x in {62,66,70} { \draw[<-] (\x,25.5) -- (\x,11.5); }
  \node[lab, fill=white, inner sep=0.6pt] at (60,18.5) {\iflangja{すべての操作}{every operation}};
  % (c) some blocks cached, frequent interaction
  \lblocks{87}{28}{1,2}{ln-accent!45}
  \draw[->, thick] (89,11.5) -- node[lab, right, pos=0.25] {\iflangja{ブロック}{blocks}} (89,25.5);
  \draw[->, thick] (109,25.5) -- node[lab, left, pos=0.25] {\iflangja{open，検証，\\書き戻し}{open, check,\\write back}} (109,11.5);
\end{tikzpicture}
   (width ~118 mm)
\begin{tikzpicture}[x=1mm, y=1mm, font=\scriptsize\sffamily,
  >={Stealth[length=1.5mm]},
  machine/.style={draw, fill=white, rounded corners=0.8mm},
  lab/.style={font=\scriptsize\sffamily, align=center},
  who/.style={font=\scriptsize\sffamily\itshape, text=ln-muted, anchor=north west},
  ttl/.style={font=\footnotesize\sffamily\bfseries, anchor=west}]
  % \lblocks{x0}{y0}{list of block indices to fill}{fill colour}
  \def\lblocks#1#2#3#4{%
    \foreach \i in {0,...,5} { \draw[fill=white] (#1+\i*4,#2) rectangle ++(4,4); }
    \foreach \i in {#3} { \draw[fill=#4] (#1+\i*4,#2) rectangle ++(4,4); }}
  \foreach \px/\t in {0/{\iflangja{(a) アップロード・ダウンロード}{(a) upload/download}},
                      40/{\iflangja{(b) 遠隔アクセス}{(b) remote access}},
                      80/{\iflangja{(c) 折衷案}{(c) compromise}}} {
    \node[ttl] at (\px,44) {\t};
    \draw[machine] (\px+1,26) rectangle (\px+37,39);
    \node[who] at (\px+1.5,38.6) {\iflangja{クライアント}{client}};
    \draw[machine, fill=ln-box] (\px+1,0) rectangle (\px+37,11);
    \node[who] at (\px+1.5,10.6) {\iflangja{サーバ}{server}};
    \lblocks{\px+7}{1.8}{0,...,5}{ln-muted!45}
  }
  % (a) whole file at the client
  \lblocks{7}{28}{0,...,5}{ln-accent!45}
  \draw[->, thick] (9,11.5) -- node[lab, right, pos=0.7] {\iflangja{ダウンロード}{download}} (9,25.5);
  \draw[->, thick] (29,25.5) -- node[lab, left, pos=0.7] {\iflangja{アップロード}{upload}} (29,11.5);
  % (b) nothing at the client, every operation to the server
  \node[lab, text=ln-muted] at (59,30) {\iflangja{ローカルにはデータなし}{no file data kept locally}};
  \foreach \x in {50,54,58} { \draw[->] (\x,25.5) -- (\x,11.5); }
  \foreach \x in {62,66,70} { \draw[<-] (\x,25.5) -- (\x,11.5); }
  \node[lab, fill=white, inner sep=0.6pt] at (60,18.5) {\iflangja{すべての操作}{every operation}};
  % (c) some blocks cached, frequent interaction
  \lblocks{87}{28}{1,2}{ln-accent!45}
  \draw[->, thick] (89,11.5) -- node[lab, right, pos=0.25] {\iflangja{ブロック}{blocks}} (89,25.5);
  \draw[->, thick] (109,25.5) -- node[lab, left, pos=0.25] {\iflangja{open，検証，\\書き戻し}{open, check,\\write back}} (109,11.5);
\end{tikzpicture}
   (width ~118 mm)
\begin{tikzpicture}[x=1mm, y=1mm, font=\scriptsize\sffamily,
  >={Stealth[length=1.5mm]},
  machine/.style={draw, fill=white, rounded corners=0.8mm},
  lab/.style={font=\scriptsize\sffamily, align=center},
  who/.style={font=\scriptsize\sffamily\itshape, text=ln-muted, anchor=north west},
  ttl/.style={font=\footnotesize\sffamily\bfseries, anchor=west}]
  % \lblocks{x0}{y0}{list of block indices to fill}{fill colour}
  \def\lblocks#1#2#3#4{%
    \foreach \i in {0,...,5} { \draw[fill=white] (#1+\i*4,#2) rectangle ++(4,4); }
    \foreach \i in {#3} { \draw[fill=#4] (#1+\i*4,#2) rectangle ++(4,4); }}
  \foreach \px/\t in {0/{\iflangja{(a) アップロード・ダウンロード}{(a) upload/download}},
                      40/{\iflangja{(b) 遠隔アクセス}{(b) remote access}},
                      80/{\iflangja{(c) 折衷案}{(c) compromise}}} {
    \node[ttl] at (\px,44) {\t};
    \draw[machine] (\px+1,26) rectangle (\px+37,39);
    \node[who] at (\px+1.5,38.6) {\iflangja{クライアント}{client}};
    \draw[machine, fill=ln-box] (\px+1,0) rectangle (\px+37,11);
    \node[who] at (\px+1.5,10.6) {\iflangja{サーバ}{server}};
    \lblocks{\px+7}{1.8}{0,...,5}{ln-muted!45}
  }
  % (a) whole file at the client
  \lblocks{7}{28}{0,...,5}{ln-accent!45}
  \draw[->, thick] (9,11.5) -- node[lab, right, pos=0.7] {\iflangja{ダウンロード}{download}} (9,25.5);
  \draw[->, thick] (29,25.5) -- node[lab, left, pos=0.7] {\iflangja{アップロード}{upload}} (29,11.5);
  % (b) nothing at the client, every operation to the server
  \node[lab, text=ln-muted] at (59,30) {\iflangja{ローカルにはデータなし}{no file data kept locally}};
  \foreach \x in {50,54,58} { \draw[->] (\x,25.5) -- (\x,11.5); }
  \foreach \x in {62,66,70} { \draw[<-] (\x,25.5) -- (\x,11.5); }
  \node[lab, fill=white, inner sep=0.6pt] at (60,18.5) {\iflangja{すべての操作}{every operation}};
  % (c) some blocks cached, frequent interaction
  \lblocks{87}{28}{1,2}{ln-accent!45}
  \draw[->, thick] (89,11.5) -- node[lab, right, pos=0.25] {\iflangja{ブロック}{blocks}} (89,25.5);
  \draw[->, thick] (109,25.5) -- node[lab, left, pos=0.25] {\iflangja{open，検証，\\書き戻し}{open, check,\\write back}} (109,11.5);
\end{tikzpicture}

  \caption{遠隔ファイルサービス．(a) クライアントはファイル全体を
    ダウンロードし，終わったらアップロードする．(b) ファイルはサーバに
    留まり，サーバがすべての操作を実行する．(c) クライアントは一部の
    ブロックをキャッシュするが，サーバと頻繁にやり取りする．}
  \label{fig:l14-remote}
\end{figure}

\begin{example}
  あるクライアントが \SI{1}{\gibi\byte} のファイルを，
  \SI[per-mode=symbol]{100}{\mebi\byte\per\second} で転送する
  ネットワーク越しに扱う．要求1回ごとに \SI{1}{\milli\second} の往復が
  かかる．\caution{ネットワークの速度はビット毎秒で表される．
  \SI{1}{\giga\bit\per\second} が運べるのは高々
  \SI{125}{\mega\byte\per\second} であり，ここで仮定した速度は
  ギガビット回線をほぼ使い切った値である．}クライアントが
  \SI{4}{\kibi\byte} のブロックを 50 個読むなら，
  遠隔アクセスモデルでは 50 回の要求，およそ \SI{50}{\milli\second} で
  済むのに対し，ファイルのダウンロードには
  $1024/100 \approx \SI{10.2}{\second}$ かかる．一方，クライアントが
  \num{200000} ブロックを読むなら，遠隔アクセスモデルでは往復だけで
  \SI{200}{\second} を要するが，ダウンロードは依然として
  \SI{10.2}{\second} であり，その後の読み出しはすべてローカルである．
  \tip{クライアントが読むブロック数が $S/(B\,t_{\mathrm{RTT}})$---1回の
  ダウンロードに収まる往復の回数---を超えると，ダウンロードが有利になる．
  ここではおよそ \num{10000} ブロックである．}
\end{example}

\subsection{折衷案}

実用的な DFS は両極端の間に位置し，次の2つの原則を組み合わせる．
\begin{enumerate}
  \item クライアントはファイルの一部，すなわち個々のブロックをローカルに
    保持してよい．ローカルなブロックに対する操作はレイテンシが低く，その分
    だけ要求が減ってサーバの負荷が下がるので，1台のサーバがより多くの
    クライアントを支えられる．
  \item クライアントはサーバとやり取りしなければならず，しかも頻繁に
    やり取りする．ファイルを open するとき，ローカルなブロックがまだ最新か
    を確かめるとき，そして自分の変更を送り返すときである．これによりサーバ
    は，クライアントが何をしているかを把握し，どのアクセスを許すかを制御
    できるようになり，一貫性の維持が容易になる．
\end{enumerate}
その代価は，サーバがより複雑になること，そしてファイル共有セマンティクスが
ローカルなファイルシステムのそれとは異なるものになることである．ブロックの
ローカルな複製の上で作業するクライアントは，他のクライアントが加えた変更を
すべて，それが起きた瞬間に観測することはできない
（\cref{sec:l14-semantics}）．

\section{ステートレスなファイルサーバとステートフルなファイルサーバ}
\label{sec:l14-state}

クライアントとサーバの間の作業の分け方は，サーバがクライアントについて
何を覚えていなければならないかを決める．\source{P4L2，ステートレスな
ファイルサーバとステートフルなファイルサーバ；OSTEP ch.~49；Silberschatz ら
ch.~19．}逆に，サーバが保持する状態は，\cref{sec:l14-remote} のモデルの
うちどれを支えられるかを決める．

\begin{definition}[ステートレスサーバ]
  \term[すてーとれすさーば]{ステートレスサーバ}[stateless server]は，
  クライアントについての情報を要求と要求の間に保持しない．どのファイルが
  open されているかも，クライアントがどのデータをキャッシュしているかも，
  クライアントがファイルのどこを読んでいるかも保持しない．すべての要求は
  自己完結的であり，それを処理するのに必要な情報をすべて含む．
\end{definition}

\begin{definition}[ステートフルサーバ]
  \term[すてーとふるさーば]{ステートフルサーバ}[stateful server]は，
  クライアントについての情報---どのクライアントがどのファイルにどの
  モードでアクセスしているか，どのデータをキャッシュあるいはロックして
  いるか---を維持する．
\end{definition}

\begin{table}[htbp]
  \centering
  \caption{ステートレスなファイルサーバとステートフルなファイルサーバ．}
  \label{tab:l14-state}
  \small
  \begin{tabular}{@{}>{\raggedright\arraybackslash}p{0.2\linewidth}>{\raggedright\arraybackslash}p{0.36\linewidth}>{\raggedright\arraybackslash}p{0.36\linewidth}@{}}
    \toprule
    & ステートレスサーバ & ステートフルサーバ \\
    \midrule
    適するモデル & アップロード・ダウンロードモデルと遠隔アクセスモデル &
      折衷モデル \\
    要求 & 自己完結的；要求当たりのビット数が多い & サーバが保持する状態を
      参照しうる \\
    支えられるもの & 一貫したキャッシュもロックも支えられない &
      ロック，キャッシュ，増分的な操作 \\
    サーバの資源 & クライアント状態のためには不要 & 状態と一貫性のために
      メモリと CPU 時間 \\
    障害の後 & 再起動して続行する & 状態をチェックポイントし復元しなければ
      ならない \\
    \bottomrule
  \end{tabular}
\end{table}

ステートレスサーバは両極端の2つのモデルには十分だが，折衷案には足りない．
クライアントがキャッシュしたブロックの一貫性を保つには，どのクライアントが
どのファイルをキャッシュしているかを知る必要があるからである．その利点は，
クライアントの状態にメモリも CPU 時間も費やさないこと，そして障害からの
回復が単に再起動するだけで済むことである．失われた状態はなく，要求に応答を
得られなかったクライアントはそれを送り直せばよい．送り直しが安全なのは，
要求が冪等である限りにおいてであり（第13章），明示的なオフセットを伴う
読み出しや書き込みはまさに冪等である．欠点は，一貫性管理を伴うキャッシュも
ロックも支えられないこと，そしてすべての要求が操作を完全に---どのファイル
か，どのオフセットか，何バイトか---記述しなければならず，要求当たりの転送
ビット数が増えることである．\caution{ステートレスとは\emph{クライアント}
についての状態を指す．サーバはファイルとそのメタデータを依然として保持して
おり，要求と要求の間に忘れるのは，誰が何を読んでいるかである．}
ステートレスサーバのクライアントでも，データを
キャッシュし，ファイルを open するときや周期的に自分の複製が最新かをサーバ
に尋ねることはできる（クライアント駆動の検証．\cref{sec:l14-caching}．
NFS がその例である）．しかしサーバは，変更をクライアントに通知することも，
複製が整合していることを保証することもできない．

折衷モデルにはステートフルサーバが必要である．サーバはクライアントが何を
しているかを知っているので，ロック，一貫性の保証を伴うクライアント
キャッシュ，そしてファイルの次のブロックを読むといった増分的な操作を
支えられる．その代価は2つある．サーバが故障すると，サーバが状態を
チェックポイントし（第8章）%
\index{ちぇっくぽいんてぃんぐ@チェックポインティング}，復旧機構がそれを
復元しない限り，状態は失われる．また，状態を維持し一貫性を強制することは
サーバの資源を消費し，その量はキャッシュ機構と一貫性プロトコルに依存する．
\cref{tab:l14-state} にこの比較をまとめる．

\begin{example}
  あるクライアントが \SI{12}{\kibi\byte} のファイルを
  \SI{4}{\kibi\byte} ずつ3回の要求で読み，2回目の読み出しの後にサーバが
  再起動する（\cref{tab:l14-state-ex}）．ステートレスサーバは読み出しの
  たびにファイルとオフセットを受け取る．ファイルハンドルが 32 バイト
  （NFSv2 での固定長），オフセットが 8 バイト，個数が 4 バイトなら，
  各要求は 44 バイトの引数を運ぶ．ステートフルサーバは open された
  ファイルとそのオフセットを記録しているので，読み出しには 4 バイトの
  ディスクリプタと個数，合わせて 8 バイトで足りる．しかし再起動の後，
  ステートフルサーバはディスクリプタ~5 をもはや知らず，状態が復元
  されない限り3回目の読み出しは失敗する．一方ステートレスサーバは，
  何事もなかったかのようにそれを処理する．
\end{example}

\begin{table}[htbp]
  \centering
  \caption{2回目の読み出しの後に再起動するステートレスサーバと
    ステートフルサーバからのファイルの読み出し．}
  \label{tab:l14-state-ex}
  \small
  \begin{tabular}{@{}l>{\raggedright\arraybackslash}p{0.38\linewidth}>{\raggedright\arraybackslash}p{0.38\linewidth}@{}}
    \toprule
    段階 & ステートレスサーバ & ステートフルサーバ \\
    \midrule
    open & \texttt{lookup(dir,"data")} がハンドル $h$ を返す &
      \texttt{open("data")} が 5 を返す；オフセット 0 を記録 \\
    read 1 & \texttt{read($h$,0,4096)} & \texttt{read(5,4096)}；オフセット
      4096 \\
    read 2 & \texttt{read($h$,4096,4096)} & \texttt{read(5,4096)}；
      オフセット 8192 \\
    再起動 & 失われるものはない & open されたファイルとオフセットが失われる \\
    read 3 & \texttt{read($h$,8192,4096)} が処理される &
      \texttt{read(5,4096)} が失敗する \\
    \bottomrule
  \end{tabular}
\end{table}

\begin{keypoint}
  実用的な DFS は，クライアントにブロックをキャッシュさせつつ，サーバとの
  やり取りを頻繁に行わせる．これはアップロード・ダウンロードモデルと
  遠隔アクセスモデルの折衷案である．この折衷にはステートフルサーバが必要
  であり，サーバの資源と障害後の復旧という代価と引き換えに，キャッシュ，
  ロック，増分的な操作を支える．ステートレスサーバはより単純で，再起動に
  よって障害を乗り切るが，クライアントのキャッシュが一貫していることを
  保証できない．
\end{keypoint}

\section{DFS におけるキャッシュ状態}
\label{sec:l14-caching}

折衷モデルでは，クライアントはファイルの状態の一部---典型的にはその
ブロックのいくつか---をローカルに保持し，open，read，write といった操作を
このキャッシュされた状態に対して行う．\source{P4L2，DFS におけるキャッシュ
状態；OSTEP ch.~49--50；Silberschatz ら ch.~19．}クライアントはこれらの
ブロックを自身のメモリに，第11章のバッファキャッシュのように保持する．
最も高速な場所だが，容量に限りがあり，クライアントがクラッシュすれば
失われる．\margin{NFS（\cref{sec:l14-nfs}）と Sprite
（\cref{sec:l14-sprite}）のクライアントは，ブロックをメモリにキャッシュ
する．}ローカルなディスクや SSD になら，クライアントはより多くのデータを
キャッシュでき，クラッシュを越えて保持できる．その代価はメモリより遅い
ことだが，それでもネットワークよりは速い．サーバもまた自身のバッファ
キャッシュにブロックをキャッシュし，それによってディスクアクセスを節約
する．どれだけ節約できるかは負荷に依存する．多数のクライアントによる
異なるファイルへの要求がそこで交錯し，互いのデータを追い出すからである．
\caution{「ネットワークより速い」は遅いネットワークを前提とする．
\num{7200}~rpm のディスクのランダムな1ブロックはおよそ
\SI{13}{\milli\second} を要し（OSTEP ch.~37），データセンタ内部の往復より
遅い．}
同じファイルが複数の場所にキャッシュされると，それらの複製は食い違い
うる．クライアント~1 と~2 がともにファイル~$F$ をキャッシュしており，
クライアント~2 が自分の複製を $F'$ に変更したとしよう
（\cref{fig:l14-caching}）．この変更はサーバに届かなければならず，
クライアント~1 は自分の $F$ の複製が古くなったことを知らなければならない．

\begin{figure}[htbp]
  \centering
  % Caching state in a DFS: client 2 modifies its cached copy of F, writes it
% back, and client 1 learns that its copy is stale, either by asking the
% server (client-driven) or by being notified (server-driven).
% Usage: % Caching state in a DFS: client 2 modifies its cached copy of F, writes it
% back, and client 1 learns that its copy is stale, either by asking the
% server (client-driven) or by being notified (server-driven).
% Usage: % Caching state in a DFS: client 2 modifies its cached copy of F, writes it
% back, and client 1 learns that its copy is stale, either by asking the
% server (client-driven) or by being notified (server-driven).
% Usage: \input{figures/l14-caching}   (width ~118 mm)
\begin{tikzpicture}[x=1mm, y=1mm, font=\scriptsize\sffamily,
  >={Stealth[length=1.5mm]},
  machine/.style={draw, rounded corners=0.8mm},
  doc/.style={draw, fill=white, minimum width=9mm, minimum height=6mm, inner sep=1pt},
  lab/.style={font=\scriptsize\sffamily, align=center},
  who/.style={font=\scriptsize\sffamily\bfseries, anchor=north west}]
  % client 1
  \draw[machine, fill=white] (0,25) rectangle (42,39);
  \node[who] at (1,38.5) {\iflangja{クライアント 1}{client 1}};
  \node[lab, anchor=west] at (3,30) {\iflangja{キャッシュ：}{cache:}};
  \node[doc] at (27,30) {$F$};
  \node[lab, text=ln-caution, anchor=west] at (32,30) {\iflangja{古い}{stale}};
  % client 2
  \draw[machine, fill=white] (76,25) rectangle (118,39);
  \node[who] at (77,38.5) {\iflangja{クライアント 2}{client 2}};
  \node[lab, anchor=west] at (79,30) {\iflangja{キャッシュ：}{cache:}};
  \node[doc, fill=ln-accent!20, thick] (f2) at (103,30) {$F'$};
  \node[lab, anchor=north, text=ln-accent] at (103,37.8) {\iflangja{(1) 書き込み}{(1) write}};
  % server
  \draw[machine, fill=ln-box] (38,0) rectangle (80,12);
  \node[who] at (39,11.5) {\iflangja{サーバ}{server}};
  \node[doc, fill=ln-accent!20] at (66,5) {$F'$};
  % (2) write back
  \draw[->, thick, ln-accent] (112,25) |- (80,4)
    node[lab, pos=0.25, left, text=ln-accent] {\iflangja{(2) 書き戻し}{(2) write back}};
  % (3a) client-driven validation
  \draw[->, densely dashed] (10,25) |- (38,4);
  \node[lab, anchor=south] at (24,4.6) {\iflangja{(3a) クライアント駆動：\\「$F$ は最新か？」}{(3a) client-driven:\\is $F$ current?}};
  % (3b) server-driven notification
  \draw[->, densely dashed, ln-caution] (52,12) -- (52,18) -| (34,25);
  \node[lab, anchor=west, text=ln-caution, align=left] at (53,17.5) {\iflangja{(3b) サーバ駆動：\\「$F$ は変更された」}{(3b) server-driven:\\$F$ has changed}};
\end{tikzpicture}
   (width ~118 mm)
\begin{tikzpicture}[x=1mm, y=1mm, font=\scriptsize\sffamily,
  >={Stealth[length=1.5mm]},
  machine/.style={draw, rounded corners=0.8mm},
  doc/.style={draw, fill=white, minimum width=9mm, minimum height=6mm, inner sep=1pt},
  lab/.style={font=\scriptsize\sffamily, align=center},
  who/.style={font=\scriptsize\sffamily\bfseries, anchor=north west}]
  % client 1
  \draw[machine, fill=white] (0,25) rectangle (42,39);
  \node[who] at (1,38.5) {\iflangja{クライアント 1}{client 1}};
  \node[lab, anchor=west] at (3,30) {\iflangja{キャッシュ：}{cache:}};
  \node[doc] at (27,30) {$F$};
  \node[lab, text=ln-caution, anchor=west] at (32,30) {\iflangja{古い}{stale}};
  % client 2
  \draw[machine, fill=white] (76,25) rectangle (118,39);
  \node[who] at (77,38.5) {\iflangja{クライアント 2}{client 2}};
  \node[lab, anchor=west] at (79,30) {\iflangja{キャッシュ：}{cache:}};
  \node[doc, fill=ln-accent!20, thick] (f2) at (103,30) {$F'$};
  \node[lab, anchor=north, text=ln-accent] at (103,37.8) {\iflangja{(1) 書き込み}{(1) write}};
  % server
  \draw[machine, fill=ln-box] (38,0) rectangle (80,12);
  \node[who] at (39,11.5) {\iflangja{サーバ}{server}};
  \node[doc, fill=ln-accent!20] at (66,5) {$F'$};
  % (2) write back
  \draw[->, thick, ln-accent] (112,25) |- (80,4)
    node[lab, pos=0.25, left, text=ln-accent] {\iflangja{(2) 書き戻し}{(2) write back}};
  % (3a) client-driven validation
  \draw[->, densely dashed] (10,25) |- (38,4);
  \node[lab, anchor=south] at (24,4.6) {\iflangja{(3a) クライアント駆動：\\「$F$ は最新か？」}{(3a) client-driven:\\is $F$ current?}};
  % (3b) server-driven notification
  \draw[->, densely dashed, ln-caution] (52,12) -- (52,18) -| (34,25);
  \node[lab, anchor=west, text=ln-caution, align=left] at (53,17.5) {\iflangja{(3b) サーバ駆動：\\「$F$ は変更された」}{(3b) server-driven:\\$F$ has changed}};
\end{tikzpicture}
   (width ~118 mm)
\begin{tikzpicture}[x=1mm, y=1mm, font=\scriptsize\sffamily,
  >={Stealth[length=1.5mm]},
  machine/.style={draw, rounded corners=0.8mm},
  doc/.style={draw, fill=white, minimum width=9mm, minimum height=6mm, inner sep=1pt},
  lab/.style={font=\scriptsize\sffamily, align=center},
  who/.style={font=\scriptsize\sffamily\bfseries, anchor=north west}]
  % client 1
  \draw[machine, fill=white] (0,25) rectangle (42,39);
  \node[who] at (1,38.5) {\iflangja{クライアント 1}{client 1}};
  \node[lab, anchor=west] at (3,30) {\iflangja{キャッシュ：}{cache:}};
  \node[doc] at (27,30) {$F$};
  \node[lab, text=ln-caution, anchor=west] at (32,30) {\iflangja{古い}{stale}};
  % client 2
  \draw[machine, fill=white] (76,25) rectangle (118,39);
  \node[who] at (77,38.5) {\iflangja{クライアント 2}{client 2}};
  \node[lab, anchor=west] at (79,30) {\iflangja{キャッシュ：}{cache:}};
  \node[doc, fill=ln-accent!20, thick] (f2) at (103,30) {$F'$};
  \node[lab, anchor=north, text=ln-accent] at (103,37.8) {\iflangja{(1) 書き込み}{(1) write}};
  % server
  \draw[machine, fill=ln-box] (38,0) rectangle (80,12);
  \node[who] at (39,11.5) {\iflangja{サーバ}{server}};
  \node[doc, fill=ln-accent!20] at (66,5) {$F'$};
  % (2) write back
  \draw[->, thick, ln-accent] (112,25) |- (80,4)
    node[lab, pos=0.25, left, text=ln-accent] {\iflangja{(2) 書き戻し}{(2) write back}};
  % (3a) client-driven validation
  \draw[->, densely dashed] (10,25) |- (38,4);
  \node[lab, anchor=south] at (24,4.6) {\iflangja{(3a) クライアント駆動：\\「$F$ は最新か？」}{(3a) client-driven:\\is $F$ current?}};
  % (3b) server-driven notification
  \draw[->, densely dashed, ln-caution] (52,12) -- (52,18) -| (34,25);
  \node[lab, anchor=west, text=ln-caution, align=left] at (53,17.5) {\iflangja{(3b) サーバ駆動：\\「$F$ は変更された」}{(3b) server-driven:\\$F$ has changed}};
\end{tikzpicture}

  \caption{クライアント 2 がキャッシュした $F$ の複製を変更し，書き戻す．
    クライアント 1 は，サーバに尋ねる（クライアント駆動）か通知を受ける
    （サーバ駆動）ことで，自分の複製が古いことを知る．}
  \label{fig:l14-caching}
\end{figure}

これは共有メモリ型マルチプロセッサにおけるキャッシュ一貫性
\index{きゃっしゅいっかんせい@キャッシュ一貫性}の問題（第10章）と同じで
ある．そこでは，ある CPU による書き込みが他の CPU のキャッシュ内の
メモリ位置の複製を古くし，キャッシュ一貫性を持つプラットフォームでは
ハードウェアがそれらを更新あるいは無効化する．DFS では一貫性の機構は
ソフトウェアで実装され，次のいずれかである．
\begin{description}
  \item[クライアント駆動] クライアントが，ファイルを open するときあるいは
    周期的に，キャッシュしたデータがまだ最新かをサーバに尋ねる．
  \item[サーバ駆動] どのクライアントがファイルをキャッシュしているかを
    知っているサーバが，ファイルが変更されたときにそれらに通知する．
    \sidenote{Sprite と同じころカーネギーメロン大学で作られた Andrew
    ファイルシステムは，サーバ駆動の設計をさらに推し進めている．
    クライアントはファイル全体を自分のローカルディスクにキャッシュし，
    サーバは，ファイルが変更される前にそれを知らせるという
    \emph{コールバック}の約束を保持する．約束が有効な間，open は何の代価も
    要さない．OSTEP ch.~50 に記述がある．}
\end{description}
クライアント駆動の確認はサーバを単純に保つが，1回ごとに要求の代価がかかり，
次の確認まではクライアントが古いデータを持ち続ける．サーバ駆動の通知は
ステートフルサーバを必要とする．いずれの機構がいつ，どれだけ頻繁に働くか
は，DFS が提供するファイル共有セマンティクスによって決まる
（\cref{sec:l14-semantics}）．

\section{ファイル共有セマンティクス}
\label{sec:l14-semantics}

単一のマシン上では，プロセス~B も open しているファイルにプロセス~A が
書き込むと，B の次の読み出しは新しいデータを返す．両プロセスとも同じ
バッファキャッシュを通してファイルにアクセスするので，すべての書き込みが
直ちに見えるからである．\source{P4L2，DFS のファイル共有セマンティクス；
Silberschatz ら ch.~15；OSTEP ch.~49--50．}クライアントがデータを
キャッシュする DFS では，クライアント~A による書き込みはまず A の
キャッシュだけを変更する．それは A がサーバへ送ったときにサーバに届き，
クライアント~B は，B がそれをサーバから取得した後にはじめてそれを見る．
あるクライアントが別のクライアントの書き込みをいつ観測するかについての
保証が，DFS の\emph{ファイル共有セマンティクス}である．

\begin{definition}[UNIXセマンティクス]
  \term[UNIXせまんてぃくす]{UNIXセマンティクス}[UNIX semantics]の下では，
  ファイルへのすべての書き込みは，どのマシン上であれ，その後のすべての
  読み出しから直ちに見える．
\end{definition}

UNIXセマンティクスは，アプリケーションがローカルなファイルシステムに期待
するものである．しかし DFS でこれを提供するには，他のどのクライアントが
ファイルを読むよりも前にすべての書き込みを伝播しなければならず，その代価は
大きい．

\begin{definition}[セッションセマンティクス]
  \term[せっしょんせまんてぃくす]{セッションセマンティクス}[session semantics]%
  の下では，クライアントはファイルを open するときにサーバから最新の版を
  取得し，ファイルを close するときに自分の変更をサーバへ書き戻す．
  open からそれに対応する close までの期間が\emph{セッション}である．
  セッションの間に他のクライアントが加えた変更は，そのセッションからは
  見えない．
\end{definition}

セッションセマンティクスは理解しやすいが，不十分なこともある．ともに
ファイルを open したままそれを共有するクライアントは互いの変更を見ず，
セッションが長く続くほど，異なるクライアントの複製が不整合なままである
時間も長くなる．

\emph{周期的な更新}では，クライアントは自分の変更を周期的にサーバへ書き戻し，
キャッシュされたデータも周期的に無効化あるいは再検証される．この周期が，
クライアントが不整合なデータを観測しうる時間の上限を与える．また DFS は，
アプリケーションが自分の変更を直ちにサーバへ反映させるための
\texttt{flush} や \texttt{sync} といった明示的な操作を提供してもよい．
こうしてクライアントがキャッシュしたデータを使う権利は時間的に制限される．
すなわちクライアントはデータに対するリースを保持している．

\begin{definition}[リース]
  \term[りーす]{リース}[lease]とは，サーバがクライアントに限られた期間
  だけ与える権利であり，キャッシュしたデータを使う権利やロックを保持する
  権利がこれにあたる．期間が満了すると，クライアントが更新していない限り
  その権利は失われる．
\end{definition}

キャッシュされたデータに対するリースは排他的である必要はない．複数の
クライアントが同じファイルに対するリースを同時に保持できる．
さらに2つの方法が，不整合な見え方そのものを回避する．

\begin{definition}[不変ファイル]
  \term[ふへんふぁいる]{不変ファイル}[immutable file]とは，作成された後は
  決して変更されないファイルである．内容を変えるには新しいファイルを作成し，
  古いものは削除してもよい．
\end{definition}

不変ファイルの複製はすべて同一なので，それを共有しても不整合は生じえない．
たとえば写真を保管するサービスは，アップロードされた写真をその場で変更する
ことは通常なく，編集された写真は新しいファイルとして保存される．
\margin{第8章のコピーオンライトと同じ仕掛けである．新しい版は別の場所に
書かれ，古いブロックは後で解放される．}最後に，
DFS は\emph{トランザクション}を提供することもできる．アプリケーションは，
1つあるいは複数のファイルに対する複数の更新を1つのトランザクションに
まとめ，DFS はそれらを，データベースシステムがそうするように，原子的に
---すべてが反映されるか，どれも反映されないかのいずれかに---適用する．

\begin{example}
  \label{ex:l14-semantics}
  ファイル $F$ の初期内容は \texttt{x} である．クライアント~A は
  $t=\SI{0}{\second}$ に $F$ を open し，\SI{4}{\second} に \texttt{y} を
  書き，\SI{10}{\second} に $F$ を close する．クライアント~B は
  \SI{2}{\second} に $F$ を open し，\SI{6}{\second} と
  \SI{12}{\second} に読み出し，\SI{14}{\second} に close する．さらに
  \SI{15}{\second} に $F$ をもう一度 open し，\SI{16}{\second} に読み出す
  （\cref{fig:l14-semantics}）．周期的な更新については，A のクライアントが
  ダーティなデータを \SI{5}{\second} ごと，すなわち 5, 10, 15,
  \ldots~秒に書き戻し，B のクライアントが 3, 8, 13, \ldots~秒に自分の
  キャッシュを再検証するものとする．B が何を読むかを
  \cref{tab:l14-semantics} に示す．セッションセマンティクスの下では，B は
  2回目のセッションではじめて \texttt{y} を見る．周期的な更新の下では，
  A の書き戻しに続く最初の再検証の後，すなわち書き込みの \SI{4}{\second}
  後に \texttt{y} を見る．一般に，どの書き込みも2つの周期の和である
  \SI{10}{\second} より長く B から見えないままになることはない．
\end{example}

\begin{figure}[htbp]
  \centering
  % File sharing semantics: client A writes y into F during a session, client
% B reads F.  (a) Session semantics: write-back on close, fetch on open.
% (b) Periodic updates: A writes back every 5 s, B revalidates at 3, 8, 13 s.
% Usage: % File sharing semantics: client A writes y into F during a session, client
% B reads F.  (a) Session semantics: write-back on close, fetch on open.
% (b) Periodic updates: A writes back every 5 s, B revalidates at 3, 8, 13 s.
% Usage: % File sharing semantics: client A writes y into F during a session, client
% B reads F.  (a) Session semantics: write-back on close, fetch on open.
% (b) Periodic updates: A writes back every 5 s, B revalidates at 3, 8, 13 s.
% Usage: \input{figures/l14-semantics}   (width ~116 mm)
\begin{tikzpicture}[x=1mm, y=1mm, font=\scriptsize\sffamily,
  >={Stealth[length=1.3mm]},
  row/.style={font=\scriptsize\sffamily, anchor=east},
  lab/.style={font=\scriptsize\sffamily, inner sep=0.8pt},
  val/.style={font=\scriptsize\ttfamily, inner sep=0.8pt},
  sess/.style={draw, fill=ln-box},
  ev/.style={circle, fill=black, inner sep=0pt, minimum size=1.3mm},
  ttl/.style={font=\footnotesize\sffamily\bfseries, anchor=west}]
  % time t (s) -> x = 20 + 5.6 t (mm)
  % \lpanel{title}  rows: A at 14, server at 7, B at 0 (inside a scope)
  \def\lpanel#1{%
    \node[ttl] at (0,21) {#1};
    \node[row] at (18,14) {\iflangja{クライアント A}{client A}};
    \node[row] at (18,7) {\iflangja{サーバ}{server}};
    \node[row] at (18,0) {\iflangja{クライアント B}{client B}};
    \foreach \y in {0,7,14} { \draw[ln-rule] (20,\y) -- (116,\y); }
    % sessions
    \draw[sess] ({20+5.6*0},12.6) rectangle ({20+5.6*10},15.4);
    \draw[sess] ({20+5.6*2},-1.4) rectangle ({20+5.6*14},1.4);
    \draw[sess] ({20+5.6*15},-1.4) rectangle ({20+5.6*17},1.4);
    % A writes y at 4 s
    \node[ev] at ({20+5.6*4},14) {};
    \node[lab, anchor=south] at ({20+5.6*4},15.6) {\iflangja{\texttt{y} を書く}{write \texttt{y}}};
    % B reads at 6, 12, 16 s
    \foreach \t in {6,12,16} { \node[ev] at ({20+5.6*\t},0) {}; }
  }
  % ---------------- (a) session semantics
  \begin{scope}[shift={(0,33)}]
    \lpanel{\iflangja{(a) セッションセマンティクス}{(a) session semantics}}
    \draw[->, thick, ln-accent] ({20+5.6*10},12.6) -- node[lab, right, text=ln-accent] {\iflangja{close で書き戻し}{write back on close}} ({20+5.6*10},7.6);
    \draw[->, thick, ln-accent] ({20+5.6*2},6.4) -- node[val, right] {x} ({20+5.6*2},1.6);
    \draw[->, thick, ln-accent] ({20+5.6*15},6.4) -- node[val, left] {y} ({20+5.6*15},1.6);
    \foreach \t/\v in {6/x,12/x,16/y} { \node[lab, anchor=north] at ({20+5.6*\t},-1.8) {\iflangja{読む：}{read }\texttt{\v}}; }
  \end{scope}
  % ---------------- (b) periodic updates
  \begin{scope}[shift={(0,0)}]
    \lpanel{\iflangja{(b) 周期的な更新（5 秒）}{(b) periodic updates (5 s)}}
    \draw[->, thick, ln-accent] ({20+5.6*5},12.6) -- node[lab, right, text=ln-accent] {\iflangja{書き戻し}{write back}} ({20+5.6*5},7.6);
    \foreach \t/\v in {3/x,8/y,13/y} {
      \draw[->, densely dashed] ({20+5.6*\t},6.4) -- node[val, right] {\v} ({20+5.6*\t},1.6);
    }
    \foreach \t/\v in {6/x,12/y,16/y} { \node[lab, anchor=north] at ({20+5.6*\t},-1.8) {\iflangja{読む：}{read }\texttt{\v}}; }
  \end{scope}
  % time axis
  \draw[->] (20,-9) -- (117,-9);
  \foreach \t in {0,5,10,15} {
    \draw ({20+5.6*\t},-8.2) -- ({20+5.6*\t},-9.8);
    \node[lab, anchor=north] at ({20+5.6*\t},-10) {\t};
  }
  \node[lab, anchor=south east] at (117,-8.6) {\iflangja{時間（秒）}{time (s)}};
\end{tikzpicture}
   (width ~116 mm)
\begin{tikzpicture}[x=1mm, y=1mm, font=\scriptsize\sffamily,
  >={Stealth[length=1.3mm]},
  row/.style={font=\scriptsize\sffamily, anchor=east},
  lab/.style={font=\scriptsize\sffamily, inner sep=0.8pt},
  val/.style={font=\scriptsize\ttfamily, inner sep=0.8pt},
  sess/.style={draw, fill=ln-box},
  ev/.style={circle, fill=black, inner sep=0pt, minimum size=1.3mm},
  ttl/.style={font=\footnotesize\sffamily\bfseries, anchor=west}]
  % time t (s) -> x = 20 + 5.6 t (mm)
  % \lpanel{title}  rows: A at 14, server at 7, B at 0 (inside a scope)
  \def\lpanel#1{%
    \node[ttl] at (0,21) {#1};
    \node[row] at (18,14) {\iflangja{クライアント A}{client A}};
    \node[row] at (18,7) {\iflangja{サーバ}{server}};
    \node[row] at (18,0) {\iflangja{クライアント B}{client B}};
    \foreach \y in {0,7,14} { \draw[ln-rule] (20,\y) -- (116,\y); }
    % sessions
    \draw[sess] ({20+5.6*0},12.6) rectangle ({20+5.6*10},15.4);
    \draw[sess] ({20+5.6*2},-1.4) rectangle ({20+5.6*14},1.4);
    \draw[sess] ({20+5.6*15},-1.4) rectangle ({20+5.6*17},1.4);
    % A writes y at 4 s
    \node[ev] at ({20+5.6*4},14) {};
    \node[lab, anchor=south] at ({20+5.6*4},15.6) {\iflangja{\texttt{y} を書く}{write \texttt{y}}};
    % B reads at 6, 12, 16 s
    \foreach \t in {6,12,16} { \node[ev] at ({20+5.6*\t},0) {}; }
  }
  % ---------------- (a) session semantics
  \begin{scope}[shift={(0,33)}]
    \lpanel{\iflangja{(a) セッションセマンティクス}{(a) session semantics}}
    \draw[->, thick, ln-accent] ({20+5.6*10},12.6) -- node[lab, right, text=ln-accent] {\iflangja{close で書き戻し}{write back on close}} ({20+5.6*10},7.6);
    \draw[->, thick, ln-accent] ({20+5.6*2},6.4) -- node[val, right] {x} ({20+5.6*2},1.6);
    \draw[->, thick, ln-accent] ({20+5.6*15},6.4) -- node[val, left] {y} ({20+5.6*15},1.6);
    \foreach \t/\v in {6/x,12/x,16/y} { \node[lab, anchor=north] at ({20+5.6*\t},-1.8) {\iflangja{読む：}{read }\texttt{\v}}; }
  \end{scope}
  % ---------------- (b) periodic updates
  \begin{scope}[shift={(0,0)}]
    \lpanel{\iflangja{(b) 周期的な更新（5 秒）}{(b) periodic updates (5 s)}}
    \draw[->, thick, ln-accent] ({20+5.6*5},12.6) -- node[lab, right, text=ln-accent] {\iflangja{書き戻し}{write back}} ({20+5.6*5},7.6);
    \foreach \t/\v in {3/x,8/y,13/y} {
      \draw[->, densely dashed] ({20+5.6*\t},6.4) -- node[val, right] {\v} ({20+5.6*\t},1.6);
    }
    \foreach \t/\v in {6/x,12/y,16/y} { \node[lab, anchor=north] at ({20+5.6*\t},-1.8) {\iflangja{読む：}{read }\texttt{\v}}; }
  \end{scope}
  % time axis
  \draw[->] (20,-9) -- (117,-9);
  \foreach \t in {0,5,10,15} {
    \draw ({20+5.6*\t},-8.2) -- ({20+5.6*\t},-9.8);
    \node[lab, anchor=north] at ({20+5.6*\t},-10) {\t};
  }
  \node[lab, anchor=south east] at (117,-8.6) {\iflangja{時間（秒）}{time (s)}};
\end{tikzpicture}
   (width ~116 mm)
\begin{tikzpicture}[x=1mm, y=1mm, font=\scriptsize\sffamily,
  >={Stealth[length=1.3mm]},
  row/.style={font=\scriptsize\sffamily, anchor=east},
  lab/.style={font=\scriptsize\sffamily, inner sep=0.8pt},
  val/.style={font=\scriptsize\ttfamily, inner sep=0.8pt},
  sess/.style={draw, fill=ln-box},
  ev/.style={circle, fill=black, inner sep=0pt, minimum size=1.3mm},
  ttl/.style={font=\footnotesize\sffamily\bfseries, anchor=west}]
  % time t (s) -> x = 20 + 5.6 t (mm)
  % \lpanel{title}  rows: A at 14, server at 7, B at 0 (inside a scope)
  \def\lpanel#1{%
    \node[ttl] at (0,21) {#1};
    \node[row] at (18,14) {\iflangja{クライアント A}{client A}};
    \node[row] at (18,7) {\iflangja{サーバ}{server}};
    \node[row] at (18,0) {\iflangja{クライアント B}{client B}};
    \foreach \y in {0,7,14} { \draw[ln-rule] (20,\y) -- (116,\y); }
    % sessions
    \draw[sess] ({20+5.6*0},12.6) rectangle ({20+5.6*10},15.4);
    \draw[sess] ({20+5.6*2},-1.4) rectangle ({20+5.6*14},1.4);
    \draw[sess] ({20+5.6*15},-1.4) rectangle ({20+5.6*17},1.4);
    % A writes y at 4 s
    \node[ev] at ({20+5.6*4},14) {};
    \node[lab, anchor=south] at ({20+5.6*4},15.6) {\iflangja{\texttt{y} を書く}{write \texttt{y}}};
    % B reads at 6, 12, 16 s
    \foreach \t in {6,12,16} { \node[ev] at ({20+5.6*\t},0) {}; }
  }
  % ---------------- (a) session semantics
  \begin{scope}[shift={(0,33)}]
    \lpanel{\iflangja{(a) セッションセマンティクス}{(a) session semantics}}
    \draw[->, thick, ln-accent] ({20+5.6*10},12.6) -- node[lab, right, text=ln-accent] {\iflangja{close で書き戻し}{write back on close}} ({20+5.6*10},7.6);
    \draw[->, thick, ln-accent] ({20+5.6*2},6.4) -- node[val, right] {x} ({20+5.6*2},1.6);
    \draw[->, thick, ln-accent] ({20+5.6*15},6.4) -- node[val, left] {y} ({20+5.6*15},1.6);
    \foreach \t/\v in {6/x,12/x,16/y} { \node[lab, anchor=north] at ({20+5.6*\t},-1.8) {\iflangja{読む：}{read }\texttt{\v}}; }
  \end{scope}
  % ---------------- (b) periodic updates
  \begin{scope}[shift={(0,0)}]
    \lpanel{\iflangja{(b) 周期的な更新（5 秒）}{(b) periodic updates (5 s)}}
    \draw[->, thick, ln-accent] ({20+5.6*5},12.6) -- node[lab, right, text=ln-accent] {\iflangja{書き戻し}{write back}} ({20+5.6*5},7.6);
    \foreach \t/\v in {3/x,8/y,13/y} {
      \draw[->, densely dashed] ({20+5.6*\t},6.4) -- node[val, right] {\v} ({20+5.6*\t},1.6);
    }
    \foreach \t/\v in {6/x,12/y,16/y} { \node[lab, anchor=north] at ({20+5.6*\t},-1.8) {\iflangja{読む：}{read }\texttt{\v}}; }
  \end{scope}
  % time axis
  \draw[->] (20,-9) -- (117,-9);
  \foreach \t in {0,5,10,15} {
    \draw ({20+5.6*\t},-8.2) -- ({20+5.6*\t},-9.8);
    \node[lab, anchor=north] at ({20+5.6*\t},-10) {\t};
  }
  \node[lab, anchor=south east] at (117,-8.6) {\iflangja{時間（秒）}{time (s)}};
\end{tikzpicture}

  \caption{\cref{ex:l14-semantics} の状況．網掛けの帯がセッションである．
    (a) セッションセマンティクスの下では，A は close で書き戻し，B は
    open でファイルを取得する．(b) 周期的な更新の下では，A は
    \SI{5}{\second} に書き戻し，B は \SI{5}{\second} ごとに再検証する
    （破線）．}
  \label{fig:l14-semantics}
\end{figure}

\begin{table}[htbp]
  \centering
  \caption{\cref{ex:l14-semantics} でクライアント B が読む $F$ の内容．}
  \label{tab:l14-semantics}
  \small
  \begin{tabular}{@{}lccc@{}}
    \toprule
    B が読む時刻 & UNIXセマンティクス & セッションセマンティクス &
      周期的な更新 \\
    \midrule
    \SI{6}{\second}  & \texttt{y} & \texttt{x} & \texttt{x} \\
    \SI{12}{\second} & \texttt{y} & \texttt{x} & \texttt{y} \\
    \SI{16}{\second} & \texttt{y} & \texttt{y} & \texttt{y} \\
    \bottomrule
  \end{tabular}
\end{table}

\needspace{6\baselineskip}
\section{ファイルとディレクトリ}
\label{sec:l14-dirs}

どの DFS にも最適といえる唯一のセマンティクスは存在しない．\source{P4L2，
ファイルサービスとディレクトリサービス．}適切な選択はワークロードに依存
する．ファイルがどれだけ頻繁に共有されるか，どれだけ頻繁に書き込まれるか，
そして一貫した見え方がアプリケーションにとってどれだけ重要かである．
こうしたアクセスパターンを知っていれば，設計者は頻繁に起こる場合の最適化
\index{ひんぱんにおこるばあいのさいてきか@頻繁に起こる場合の最適化}
（第1章）を行える．

DFS は，ワークロードの異なる2種類のファイルを扱う．通常のファイルは
たいてい一度に1つのクライアントだけが使い，変更が他のクライアントから
見えるまでに遅れがあっても気づかれることはまれである．ディレクトリは
はるかに頻繁に共有される---共有ディレクトリでファイルを作成，削除，
改名するクライアントはすべてそれを変更する---うえ，ディレクトリの
不整合な見え方は有害である．たとえば，2つのクライアントが同じ名前の
ファイルをともに作成してしまいかねない．そのため DFS は，両者に異なる
方針を適用できる．たとえば，ファイルにはセッションセマンティクス，
ディレクトリには UNIXセマンティクスを用いる，あるいは周期的な更新を
用いつつ，書き戻しの頻度をファイルについてはディレクトリより低くする，
といった具合である．

\begin{keypoint}
  ファイル共有セマンティクスは，あるクライアントが別のクライアントの
  書き込みをいつ観測するかを定める．UNIXセマンティクスはすべての書き込みを
  直ちに見えるようにするが DFS では高価であり，セッションセマンティクスは
  close と open で変更を伝播し，周期的な更新は不整合を更新周期で抑える．
  不変ファイルとトランザクションは不整合な見え方そのものを回避する．
  選択はワークロードに従い，ファイルとディレクトリを別に扱ってもよい．
\end{keypoint}

\section{レプリケーションとパーティショニング}
\label{sec:l14-replication}

複数のマシン上で動作するファイルサービスは，ファイルをそれらのマシンに
2通りのいずれかの方法で分散させる（\cref{fig:l14-models}b と~c）．
\source{P4L2，レプリケーションとパーティショニング；Silberschatz ら
ch.~19．}

\begin{definition}[レプリケーション]
  \term[れぷりけーしょん]{レプリケーション}[replication]では，どのサーバ
  マシンも，すべてのファイルの完全な複製，すなわち\emph{レプリカ}を保持
  する．
\end{definition}

レプリケーションは負荷を均す．読み出しはどのレプリカでも処理できるから
である．また可用性と耐障害性をもたらす．1台のマシンが故障しても，残りが
すべてのファイルを提供し続けられる．しかし書き込みはより複雑になる．
すべての書き込みをすべてのレプリカに対して同期的に行うか---この場合
書き込みは遅くなる---，あるいは1つのレプリカに対して行って後から他へ
伝播するかであり，後者ではしばらくの間レプリカが食い違う．食い違った
レプリカは整合させなければならない．たとえば投票によって，過半数の
レプリカが保持する内容を採用する．

\begin{definition}[パーティショニング]
  \term[ぱーてぃしょにんぐ]{パーティショニング}[partitioning]では，各サーバ
  マシンがファイルの部分集合を保持し，それらの部分集合が合わさって
  ファイルシステムを構成する．
\end{definition}

パーティショニングは単一のサーバに比べて可用性を高める．1台のマシンの故障は
そのマシンが保持するファイルにしか影響しないからである．またファイル
システムの規模に対してスケールする．ファイルが増えてもマシンを増やせば
よい．1つのファイルへの書き込みにはそれを保持するマシンだけが関わるので，
書き込みは単純なままである．他方で，故障したマシンが保持していたデータは
失われ，負荷の均衡をとるのはより難しい．人気のあるファイルが少数のマシンに
集中すると，それらがホットスポットになる．\cref{tab:l14-repl} に両者を
比較する．2つを組み合わせることもできる．ファイルを分割し，各パーティション
を複数のマシンに複製するのである．%
\begin{marginfigure}
  \centering
  % Replication combined with partitioning: two partitions of 250 files, each
% replicated on two of four machines.
% Usage: % Replication combined with partitioning: two partitions of 250 files, each
% replicated on two of four machines.
% Usage: % Replication combined with partitioning: two partitions of 250 files, each
% replicated on two of four machines.
% Usage: \input{figures/l14-combined}   (margin figure, width ~48 mm)
\begin{tikzpicture}[x=1mm, y=1mm, font=\scriptsize\sffamily,
  sv/.style={draw, fill=ln-box, minimum width=18mm, minimum height=8mm, inner sep=1pt, align=center},
  grp/.style={draw, densely dashed, ln-muted, rounded corners=0.8mm},
  lab/.style={font=\scriptsize\sffamily, text=ln-muted, anchor=south west, inner sep=0.5pt}]
  \foreach \y/\p/\f/\a/\b in {16/1/1--250/1/2, 0/2/251--500/3/4} {
    \draw[grp] (0,\y-5) rectangle (47,\y+9);
    \node[lab] at (1,\y+5) {\iflangja{パーティション}{partition} \p: \iflangja{ファイル}{files} \f};
    \node[sv] at (11.5,\y) {S\a};
    \node[sv] at (35.5,\y) {S\b};
  }
\end{tikzpicture}
   (margin figure, width ~48 mm)
\begin{tikzpicture}[x=1mm, y=1mm, font=\scriptsize\sffamily,
  sv/.style={draw, fill=ln-box, minimum width=18mm, minimum height=8mm, inner sep=1pt, align=center},
  grp/.style={draw, densely dashed, ln-muted, rounded corners=0.8mm},
  lab/.style={font=\scriptsize\sffamily, text=ln-muted, anchor=south west, inner sep=0.5pt}]
  \foreach \y/\p/\f/\a/\b in {16/1/1--250/1/2, 0/2/251--500/3/4} {
    \draw[grp] (0,\y-5) rectangle (47,\y+9);
    \node[lab] at (1,\y+5) {\iflangja{パーティション}{partition} \p: \iflangja{ファイル}{files} \f};
    \node[sv] at (11.5,\y) {S\a};
    \node[sv] at (35.5,\y) {S\b};
  }
\end{tikzpicture}
   (margin figure, width ~48 mm)
\begin{tikzpicture}[x=1mm, y=1mm, font=\scriptsize\sffamily,
  sv/.style={draw, fill=ln-box, minimum width=18mm, minimum height=8mm, inner sep=1pt, align=center},
  grp/.style={draw, densely dashed, ln-muted, rounded corners=0.8mm},
  lab/.style={font=\scriptsize\sffamily, text=ln-muted, anchor=south west, inner sep=0.5pt}]
  \foreach \y/\p/\f/\a/\b in {16/1/1--250/1/2, 0/2/251--500/3/4} {
    \draw[grp] (0,\y-5) rectangle (47,\y+9);
    \node[lab] at (1,\y+5) {\iflangja{パーティション}{partition} \p: \iflangja{ファイル}{files} \f};
    \node[sv] at (11.5,\y) {S\a};
    \node[sv] at (35.5,\y) {S\b};
  }
\end{tikzpicture}

  \caption{2つのパーティションを，それぞれ2台のマシンに複製した構成．}
  \label{fig:l14-combined}
\end{marginfigure}

\begin{table}[htbp]
  \centering
  \caption{レプリケーションとパーティショニング．}
  \label{tab:l14-repl}
  \small
  \begin{tabular}{@{}>{\raggedright\arraybackslash}p{0.2\linewidth}>{\raggedright\arraybackslash}p{0.36\linewidth}>{\raggedright\arraybackslash}p{0.36\linewidth}@{}}
    \toprule
    & レプリケーション & パーティショニング \\
    \midrule
    容量 & 1台のマシンの容量 & 全マシンの合計 \\
    読み出し & どのレプリカでも処理できる；負荷が均される &
      そのファイルを保持するマシンが処理する；ホットスポットが生じうる \\
    書き込み & すべてのレプリカに対して，同期的にあるいは後から伝播；
      レプリカを整合させる必要がある & 1台のマシン上で；単純 \\
    マシンの故障 & ファイルは失われない & そのマシンのファイルが失われる \\
    \bottomrule
  \end{tabular}
\end{table}

\begin{example}
  4台のサーバマシンがそれぞれ 250 個のファイルを格納できるとする．複製
  すれば DFS は 250 個のファイルを保持し，1台の---実際には3台までの---
  マシンの故障でも1つも失われない．分割すれば \num{1000} 個のファイルを
  保持するが，1台の故障でそのうち $250/1000 = \SI{25}{\percent}$ が
  失われる．\cref{fig:l14-combined} のように組み合わせ，250 個ずつの
  2つのパーティションをそれぞれ2台のマシンに複製すると，DFS は 500 個の
  ファイルを保持する．1台の故障では1つも失われず，同じパーティションの
  2台がともに故障したときにのみ \SI{50}{\percent} が失われる．
\end{example}

\section{Network File System}
\label{sec:l14-nfs}

1980 年代に Sun Microsystems で開発された \term{NFS}[Network File System]%
\index{Network File System|see{NFS}}（Network File System）\cite{sandberg1985}
は，クライアントがリモートのサーバ上にあるファイルにアクセスする，広く
使われている DFS である．\source{P4L2，NFS の設計；Sandberg ら
\cite{sandberg1985}；OSTEP ch.~49；Silberschatz ら ch.~15．}その構造を
\cref{fig:l14-nfs} に示す．クライアント上のアプリケーションは通常の
ファイルシステムコールを発行する．VFS は，そのファイルが NFS ファイル
システムに属することを見出し，操作を\emph{NFS クライアント}へ渡す．これは，
操作をサーバへの遠隔手続き呼び出しに変換するファイルシステム実装である．
サーバ側では\emph{NFS サーバ}の構成要素が要求を受け取り，サーバ自身の VFS と
ローカルファイルシステムを通じてサーバ上のファイルに対して操作を行い，
結果を返す．結果はアプリケーションまで戻っていく．NFS は Sun RPC
\index{Sun RPC} とそのデータ表現 XDR（第13章）の上に構築されている．

\begin{figure}[htbp]
  \centering
  % NFS architecture: the VFS on the client passes operations on NFS files to
% the NFS client, which sends RPC requests to the NFS server; the server
% performs them through its own VFS and local file system.
% Usage: % NFS architecture: the VFS on the client passes operations on NFS files to
% the NFS client, which sends RPC requests to the NFS server; the server
% performs them through its own VFS and local file system.
% Usage: % NFS architecture: the VFS on the client passes operations on NFS files to
% the NFS client, which sends RPC requests to the NFS server; the server
% performs them through its own VFS and local file system.
% Usage: \input{figures/l14-nfs}   (width ~118 mm)
\begin{tikzpicture}[x=1mm, y=1mm, font=\scriptsize\sffamily,
  >={Stealth[length=1.5mm]},
  machine/.style={draw, rounded corners=0.8mm, fill=white},
  box/.style={draw, fill=white, minimum height=6mm, inner sep=1pt, align=center},
  kbox/.style={draw, fill=ln-box, minimum height=6mm, inner sep=1pt, align=center},
  nfs/.style={draw, fill=ln-accent!18, minimum height=6mm, inner sep=1pt, align=center},
  dev/.style={draw, fill=ln-accent!35, minimum height=5mm, inner sep=1pt, align=center},
  lab/.style={font=\scriptsize\sffamily, align=center},
  who/.style={font=\footnotesize\sffamily\bfseries, anchor=north west}]
  % ---------------- client machine
  \draw[machine] (0,0) rectangle (56,68);
  \node[who] at (1,67.5) {\iflangja{クライアントマシン}{client machine}};
  \node[box, minimum width=44mm] (app) at (28,56) {\iflangja{アプリケーション}{applications}};
  \draw[densely dashed, ln-muted] (1,50.5) -- (55,50.5);
  \node[kbox, minimum width=44mm] (sc) at (28,45) {\iflangja{システムコール層}{system call layer}};
  \node[kbox, minimum width=44mm] (vfs) at (28,35) {VFS};
  \node[kbox, minimum width=20mm] (lfs) at (16,24) {\iflangja{ローカル FS}{local FS}};
  \node[nfs, minimum width=20mm] (nc) at (40,24) {\iflangja{NFS クライアント}{NFS client}};
  \node[dev, minimum width=20mm] (d1) at (16,12) {\iflangja{ディスク}{disk}};
  \node[nfs, minimum width=20mm] (cs) at (40,12) {\iflangja{RPC スタブ}{RPC stub}};
  \draw[<->, thick] (app) -- (sc);
  \draw[<->, thick] (sc) -- (vfs);
  \draw[<->, thick] (lfs.north |- vfs.south) -- (lfs.north);
  \draw[<->, thick] (nc.north |- vfs.south) -- (nc.north);
  \draw[<->, thick] (lfs) -- (d1);
  \draw[<->, thick] (nc) -- (cs);
  % ---------------- server machine
  \draw[machine] (62,0) rectangle (118,68);
  \node[who] at (63,67.5) {\iflangja{サーバマシン}{server machine}};
  \node[nfs, minimum width=44mm] (ss) at (90,56) {\iflangja{RPC スタブ}{RPC stub}};
  \node[nfs, minimum width=44mm] (ns) at (90,45) {\iflangja{NFS サーバ}{NFS server}};
  \node[kbox, minimum width=44mm] (svfs) at (90,35) {VFS};
  \node[kbox, minimum width=44mm] (slfs) at (90,24) {\iflangja{ローカル FS}{local FS}};
  \node[dev, minimum width=44mm] (d2) at (90,12) {\iflangja{ディスク}{disk}};
  \draw[<->, thick] (ss) -- (ns);
  \draw[<->, thick] (ns) -- (svfs);
  \draw[<->, thick] (svfs) -- (slfs);
  \draw[<->, thick] (slfs) -- (d2);
  % ---------------- network
  \draw[<->, line width=1.4pt, ln-accent] (cs.east) -- (59,12) |- (ss.west);
  \node[lab, rotate=90, fill=white, inner sep=0.8pt, text=ln-accent] at (59,34) {\iflangja{ネットワーク（RPC）}{network (RPC)}};
  \node[lab, text=ln-muted, align=left, anchor=west] at (66,4.2)
    {\iflangja{要求はファイルハンドルで対象ファイルを指定}{requests name files by their file handles}};
\end{tikzpicture}
   (width ~118 mm)
\begin{tikzpicture}[x=1mm, y=1mm, font=\scriptsize\sffamily,
  >={Stealth[length=1.5mm]},
  machine/.style={draw, rounded corners=0.8mm, fill=white},
  box/.style={draw, fill=white, minimum height=6mm, inner sep=1pt, align=center},
  kbox/.style={draw, fill=ln-box, minimum height=6mm, inner sep=1pt, align=center},
  nfs/.style={draw, fill=ln-accent!18, minimum height=6mm, inner sep=1pt, align=center},
  dev/.style={draw, fill=ln-accent!35, minimum height=5mm, inner sep=1pt, align=center},
  lab/.style={font=\scriptsize\sffamily, align=center},
  who/.style={font=\footnotesize\sffamily\bfseries, anchor=north west}]
  % ---------------- client machine
  \draw[machine] (0,0) rectangle (56,68);
  \node[who] at (1,67.5) {\iflangja{クライアントマシン}{client machine}};
  \node[box, minimum width=44mm] (app) at (28,56) {\iflangja{アプリケーション}{applications}};
  \draw[densely dashed, ln-muted] (1,50.5) -- (55,50.5);
  \node[kbox, minimum width=44mm] (sc) at (28,45) {\iflangja{システムコール層}{system call layer}};
  \node[kbox, minimum width=44mm] (vfs) at (28,35) {VFS};
  \node[kbox, minimum width=20mm] (lfs) at (16,24) {\iflangja{ローカル FS}{local FS}};
  \node[nfs, minimum width=20mm] (nc) at (40,24) {\iflangja{NFS クライアント}{NFS client}};
  \node[dev, minimum width=20mm] (d1) at (16,12) {\iflangja{ディスク}{disk}};
  \node[nfs, minimum width=20mm] (cs) at (40,12) {\iflangja{RPC スタブ}{RPC stub}};
  \draw[<->, thick] (app) -- (sc);
  \draw[<->, thick] (sc) -- (vfs);
  \draw[<->, thick] (lfs.north |- vfs.south) -- (lfs.north);
  \draw[<->, thick] (nc.north |- vfs.south) -- (nc.north);
  \draw[<->, thick] (lfs) -- (d1);
  \draw[<->, thick] (nc) -- (cs);
  % ---------------- server machine
  \draw[machine] (62,0) rectangle (118,68);
  \node[who] at (63,67.5) {\iflangja{サーバマシン}{server machine}};
  \node[nfs, minimum width=44mm] (ss) at (90,56) {\iflangja{RPC スタブ}{RPC stub}};
  \node[nfs, minimum width=44mm] (ns) at (90,45) {\iflangja{NFS サーバ}{NFS server}};
  \node[kbox, minimum width=44mm] (svfs) at (90,35) {VFS};
  \node[kbox, minimum width=44mm] (slfs) at (90,24) {\iflangja{ローカル FS}{local FS}};
  \node[dev, minimum width=44mm] (d2) at (90,12) {\iflangja{ディスク}{disk}};
  \draw[<->, thick] (ss) -- (ns);
  \draw[<->, thick] (ns) -- (svfs);
  \draw[<->, thick] (svfs) -- (slfs);
  \draw[<->, thick] (slfs) -- (d2);
  % ---------------- network
  \draw[<->, line width=1.4pt, ln-accent] (cs.east) -- (59,12) |- (ss.west);
  \node[lab, rotate=90, fill=white, inner sep=0.8pt, text=ln-accent] at (59,34) {\iflangja{ネットワーク（RPC）}{network (RPC)}};
  \node[lab, text=ln-muted, align=left, anchor=west] at (66,4.2)
    {\iflangja{要求はファイルハンドルで対象ファイルを指定}{requests name files by their file handles}};
\end{tikzpicture}
   (width ~118 mm)
\begin{tikzpicture}[x=1mm, y=1mm, font=\scriptsize\sffamily,
  >={Stealth[length=1.5mm]},
  machine/.style={draw, rounded corners=0.8mm, fill=white},
  box/.style={draw, fill=white, minimum height=6mm, inner sep=1pt, align=center},
  kbox/.style={draw, fill=ln-box, minimum height=6mm, inner sep=1pt, align=center},
  nfs/.style={draw, fill=ln-accent!18, minimum height=6mm, inner sep=1pt, align=center},
  dev/.style={draw, fill=ln-accent!35, minimum height=5mm, inner sep=1pt, align=center},
  lab/.style={font=\scriptsize\sffamily, align=center},
  who/.style={font=\footnotesize\sffamily\bfseries, anchor=north west}]
  % ---------------- client machine
  \draw[machine] (0,0) rectangle (56,68);
  \node[who] at (1,67.5) {\iflangja{クライアントマシン}{client machine}};
  \node[box, minimum width=44mm] (app) at (28,56) {\iflangja{アプリケーション}{applications}};
  \draw[densely dashed, ln-muted] (1,50.5) -- (55,50.5);
  \node[kbox, minimum width=44mm] (sc) at (28,45) {\iflangja{システムコール層}{system call layer}};
  \node[kbox, minimum width=44mm] (vfs) at (28,35) {VFS};
  \node[kbox, minimum width=20mm] (lfs) at (16,24) {\iflangja{ローカル FS}{local FS}};
  \node[nfs, minimum width=20mm] (nc) at (40,24) {\iflangja{NFS クライアント}{NFS client}};
  \node[dev, minimum width=20mm] (d1) at (16,12) {\iflangja{ディスク}{disk}};
  \node[nfs, minimum width=20mm] (cs) at (40,12) {\iflangja{RPC スタブ}{RPC stub}};
  \draw[<->, thick] (app) -- (sc);
  \draw[<->, thick] (sc) -- (vfs);
  \draw[<->, thick] (lfs.north |- vfs.south) -- (lfs.north);
  \draw[<->, thick] (nc.north |- vfs.south) -- (nc.north);
  \draw[<->, thick] (lfs) -- (d1);
  \draw[<->, thick] (nc) -- (cs);
  % ---------------- server machine
  \draw[machine] (62,0) rectangle (118,68);
  \node[who] at (63,67.5) {\iflangja{サーバマシン}{server machine}};
  \node[nfs, minimum width=44mm] (ss) at (90,56) {\iflangja{RPC スタブ}{RPC stub}};
  \node[nfs, minimum width=44mm] (ns) at (90,45) {\iflangja{NFS サーバ}{NFS server}};
  \node[kbox, minimum width=44mm] (svfs) at (90,35) {VFS};
  \node[kbox, minimum width=44mm] (slfs) at (90,24) {\iflangja{ローカル FS}{local FS}};
  \node[dev, minimum width=44mm] (d2) at (90,12) {\iflangja{ディスク}{disk}};
  \draw[<->, thick] (ss) -- (ns);
  \draw[<->, thick] (ns) -- (svfs);
  \draw[<->, thick] (svfs) -- (slfs);
  \draw[<->, thick] (slfs) -- (d2);
  % ---------------- network
  \draw[<->, line width=1.4pt, ln-accent] (cs.east) -- (59,12) |- (ss.west);
  \node[lab, rotate=90, fill=white, inner sep=0.8pt, text=ln-accent] at (59,34) {\iflangja{ネットワーク（RPC）}{network (RPC)}};
  \node[lab, text=ln-muted, align=left, anchor=west] at (66,4.2)
    {\iflangja{要求はファイルハンドルで対象ファイルを指定}{requests name files by their file handles}};
\end{tikzpicture}

  \caption{NFS．クライアントの VFS は NFS ファイルに対する操作を NFS
    クライアントへ渡し，NFS クライアントはそれを RPC 要求として NFS
    サーバへ送る．サーバは自身の VFS とローカルファイルシステムを通じて
    それを実行する．}
  \label{fig:l14-nfs}
\end{figure}

クライアントがどのリモートファイルに到達できるかは，NFS ファイルシステムが
クライアントのディレクトリツリーに取り付けられるときに決まる．サーバは
自分のディレクトリのいくつかを\emph{エクスポート}し，クライアントは
\term[まうんと]{マウント}[mount]操作でエクスポートされたディレクトリを
利用可能にする．マウントはそれをローカルなディレクトリ，すなわち
\emph{マウントポイント}に取り付ける．マウントポイント以下のパスはサーバ上の
ファイルを指すことになる．プロセスがそのようなファイルを open すると，
NFS クライアントはそのパスをサーバで解決し，サーバはファイルハンドルを
返す．

\begin{definition}[ファイルハンドル]
  \term[ふぁいるはんどる]{ファイルハンドル}[file handle]とは，NFS
  クライアントと NFS サーバが要求の中でファイルを識別するために用いる
  バイト列である．これを作るのはサーバであり，ファイルのサーバローカルな
  同一性---ファイルシステムと，その中でのファイル---を符号化する．
  クライアントはそれを不透明なものとして扱う．
\end{definition}

NFS クライアントは，open されたファイルを表す自身のデータ構造にハンドルを
保持し，読み出し，書き込み，ファイルの属性の要求といったその後のすべての
要求でそれを渡す．\margin{ハンドルは NFSv2 では \SI{32}{\byte} の固定長で
ある．以降の版では可変長になり，NFSv3 では高々 \SI{64}{\byte}
\cite{rfc1813}，NFSv4 では \SI{128}{\byte} \cite{rfc7530} である．}
サーバは通常，ハンドルにファイルシステムの識別子，
ファイルの iノード番号，そして iノードが再利用されるたびに変わる\emph{世代番号}を
符号化する．ファイルがサーバ上で，ローカルなプロセスによってあるいは別の
クライアントによって削除されると，そのハンドルを用いた要求は
\emph{stale file handle}（古くなったファイルハンドル）というエラーで
失敗する．ハンドルはもはや存在しないファイルを指しており，世代番号は，
その iノードを再利用した新しいファイルをハンドルが黙って指してしまうことを
防ぐ．

\begin{example}
  あるクライアントのディレクトリ \texttt{/home} がサーバからマウントされて
  おり，マウント操作によってクライアントはエクスポートされたディレクトリの
  ファイルハンドルを得ている．次のプログラムはそれを知らない．コメントは，
  何もキャッシュされていないときに各呼び出しが引き起こす主要な NFSv3 の
  要求を示す．
  \needspace{10\baselineskip}
  \begin{lstlisting}[language=C, numbers=none]
char buf[4096];
int fd = open("/home/ana/log.txt", O_RDONLY);
        /* /home のハンドルで "ana" を LOOKUP，
           ana のハンドルで "log.txt" を LOOKUP：
           サーバがファイルハンドルを返す         */
ssize_t n = read(fd, buf, sizeof buf);
        /* READ(ファイルハンドル, オフセット 0, 4096) */
close(fd);  /* 読み出し専用：要求は不要 */
\end{lstlisting}
\end{example}

\section{NFS のバージョン}
\label{sec:l14-nfs-versions}

NFS は 1980 年代から使われており，いくつかのバージョンを経てきた．現在
使われているのは NFSv3 \cite{rfc1813} と NFSv4 \cite{rfc7530} である．
\source{P4L2，NFS のバージョン；RFC~1813 \cite{rfc1813}；RFC~7530
\cite{rfc7530}；OSTEP ch.~49．}両者は状態の扱いにおいて根本的に異なる．
NFSv3 はステートレスなプロトコルである．ただし実際には，ロックマネージャ
のような状態を保持する補助プロトコルによって補われている．\sidenote{NFSv3
は実のところ RPC プログラムの一族である．NFS 本体 \cite{rfc1813}，パスを
エクスポートされたディレクトリのハンドルに変える MOUNT プロトコル，
ロックマネージャ，そして再起動したことを互いに伝える状態モニタである．
いずれも \texttt{rpcbind}（第13章）に登録される．NFSv4 はこれらを1つの
プロトコルにまとめ，ポート 2049 だけを使う．}NFSv4 は
ステートフルであり，ファイルの open，ロック，そして後述する委譲がプロトコル
の一部になっている．

\subsection{キャッシュ}

並行にアクセスされないファイルについて，NFS はセッションに基づくセマン
ティクスを提供する．クライアントはファイルを close するときに変更をサーバへ
書き戻し，open するときに新しい版がないかを確認する．さらに，ファイルの
使用中も，クライアントは周期的に，ファイルが変更されていないかをサーバに
尋ねる．既定では，ファイルについては \SI{3}{\second} ごと，ディレクトリ
については \SI{30}{\second} ごとである．\margin{Linux では，この周期は
マウントオプション \texttt{acregmin} と \texttt{acdirmin} である．既定で
いずれも \SI{60}{\second} の \texttt{acregmax} と \texttt{acdirmax} が
上限を与える．}
ディレクトリの周期が長いのは意図的である．ディレクトリの内容はファイルの
データよりはるかに変更が少ない一方でその中の探索は頻繁なので，確認を
まれにすれば多くの要求を節約でき，代価は，別のクライアントが作成あるいは
削除した名前を知るのが遅れることだけである．この方針は
\cref{sec:l14-dirs} の例とは逆だが，それが許容されるのは，ディレクトリを
変更する操作はいずれにせよサーバへ行き，そこで直列化されるからである．

\begin{example}
  クライアント~B が最後にファイル $F$ をサーバと照合したのが
  $t=\SI{20.0}{\second}$ なら，$F$ についての B の情報は
  \SI{23.0}{\second} まで最新とみなされる．クライアント~A が
  \SI{21.0}{\second} に $F$ を close し，変更を書き戻す．B による
  \SI{22.5}{\second} の読み出しは B のキャッシュから処理され，古いデータを
  返す．\SI{23.0}{\second} より後の B の最初の読み出しで B はサーバと照合し，
  変更を見つけて新しいデータを取得する．B が \SI{20.0}{\second} に照合した
  ディレクトリに A が \SI{21.0}{\second} に作成したファイルは，
  \SI{50.0}{\second} まで B から見えないままでありうる．
\end{example}

NFSv4 はこうした確認を避けられる．サーバはファイルを一定期間クライアントに
\emph{委譲}してよい．委譲が続く間，クライアントはファイルを open し，
更新の有無をサーバに確認することなくキャッシュしたデータを使う．そして
サーバは，別のクライアントがそのファイルに対して委譲と衝突するアクセスを
必要としたときに委譲を回収する．

\begin{note}
  このように NFS は，セッションセマンティクスと周期的な再検証を組み合わせて
  いる．\margin{セッションに似た側は \emph{close-to-open} 一貫性と呼ばれる．
  open で再検証し，close で書き戻す．\texttt{nocto} マウントオプションは
  open での確認を止める．\texttt{nfs(5)}．}周期的な確認によって，まだ
  open しているセッションの中に他の
  クライアントの変更が持ち込まれることがある．純粋なセッションセマンティクス
  ではこれは起こらない．また2回の確認の間に，クライアントはサーバ上ですでに
  変更されたデータを読むことがある．したがって UNIXセマンティクスは提供
  されない．
\end{note}

\subsection{ロック}

NFSv4 のロックはリースに基づく．\margin{NFSv3 のロックマネージャは独立した
プロトコルであり，Network Lock Manager（NLM）と呼ばれる．}サーバはロックを
限られた期間だけ与え，クライアントはこのリースが満了する前に，ロックを解放
するかリースを更新しなければならない．クライアントが故障したり到達不能に
なったりすれば，そのリースは期限切れとなり，サーバは故障したクライアントを
いつまでも待つことなく，他のクライアントにロックを与えられる
（\cref{fig:l14-lease}）．NFSv4 はさらに\emph{共有予約}を支える．
クライアントはファイルを open するとき，自分が必要とするアクセス---読み
出し，書き込み，あるいはその両方---と，他のクライアントに拒否するアクセスを
宣言する．読み出しは共有しつつ他への書き込みを拒否する予約は，
リーダ・ライタロック\index{りーだらいたろっく@リーダ・ライタロック}
（第10章）のように働く．クライアントは，ファイルを close することなく，
予約を読み出しから書き込みへ昇格させたり，逆に降格させたりできる．

\begin{figure}[htbp]
  \centering
  % Lease-based locking: client 1 acquires a lock with a 30 s lease, renews it
% at 25 s and crashes at 40 s; the lease expires at 55 s and the server
% grants the lock to client 2, which has been waiting since 45 s.
% Usage: % Lease-based locking: client 1 acquires a lock with a 30 s lease, renews it
% at 25 s and crashes at 40 s; the lease expires at 55 s and the server
% grants the lock to client 2, which has been waiting since 45 s.
% Usage: % Lease-based locking: client 1 acquires a lock with a 30 s lease, renews it
% at 25 s and crashes at 40 s; the lease expires at 55 s and the server
% grants the lock to client 2, which has been waiting since 45 s.
% Usage: \input{figures/l14-lease}   (width ~116 mm)
\begin{tikzpicture}[x=1mm, y=1mm, font=\scriptsize\sffamily,
  >={Stealth[length=1.3mm]},
  row/.style={font=\scriptsize\sffamily, anchor=east},
  lab/.style={font=\scriptsize\sffamily, inner sep=0.8pt, align=center}]
  % time t (s) -> x = 18 + 1.4 t (mm)
  \node[row] at (16,16) {\iflangja{クライアント 1}{client 1}};
  \node[row] at (16,8) {\iflangja{サーバ}{server}};
  \node[row] at (16,0) {\iflangja{クライアント 2}{client 2}};
  \draw[ln-rule] (18,16) -- ({18+1.4*40},16);
  \draw[ln-rule, densely dotted] ({18+1.4*40},16) -- (116,16);
  \draw[ln-rule] (18,8) -- (116,8);
  \draw[ln-rule] (18,0) -- (116,0);
  % lock holder on the server
  \draw[fill=ln-box] ({18+1.4*0},6) rectangle ({18+1.4*55},10);
  \node[lab] at ({18+1.4*12},8) {\iflangja{クライアント 1 がロックを保持}{held by client 1}};
  \draw[fill=ln-accent!30] ({18+1.4*55},6) rectangle (116,10);
  \node[lab] at ({(18+1.4*55+116)/2},8) {\iflangja{クライアント 2}{client 2}};
  % lease expiry marks
  \foreach \t in {30,55} {
    \draw[densely dashed, ln-caution] ({18+1.4*\t},4.6) -- ({18+1.4*\t},11.4);
  }
  \node[lab, text=ln-caution, anchor=south] at ({18+1.4*55},11.4) {\iflangja{リース期限切れ}{lease expires}};
  \node[lab, text=ln-caution, anchor=north] at ({18+1.4*30},4.6) {\iflangja{当初の期限}{first expiry}};
  % client 1: acquire, renew, crash
  \draw[->, thick] ({18+1.4*0},14.8) -- ({18+1.4*0},10.3);
  \node[lab, anchor=south west] at ({18+1.4*0},17) {\iflangja{ロック獲得（リース 30 秒）}{lock (30 s lease)}};
  \draw[->, thick] ({18+1.4*25},14.8) -- ({18+1.4*25},10.3);
  \node[lab, anchor=south] at ({18+1.4*25},17) {\iflangja{更新}{renew}};
  \node[text=ln-caution, font=\normalsize\bfseries] at ({18+1.4*40},16) {$\times$};
  \node[lab, anchor=south, text=ln-caution] at ({18+1.4*40},17) {\iflangja{クラッシュ}{crash}};
  % client 2: request and grant
  \draw[->, thick] ({18+1.4*45},1.2) -- ({18+1.4*45},5.7);
  \node[lab, anchor=north] at ({18+1.4*45},-1.2) {\iflangja{ロック要求}{request}};
  \draw[densely dotted, thick] ({18+1.4*45},0) -- ({18+1.4*55},0);
  \draw[->, thick] ({18+1.4*55},5.7) -- ({18+1.4*55},1.2);
  \node[lab, anchor=north] at ({18+1.4*55},-1.2) {\iflangja{許可}{granted}};
  % time axis
  \draw[->] (18,-8) -- (117,-8);
  \foreach \t in {0,10,...,70} {
    \draw ({18+1.4*\t},-7.2) -- ({18+1.4*\t},-8.8);
    \node[lab, anchor=north] at ({18+1.4*\t},-9) {\t};
  }
  \node[lab, anchor=south east] at (117,-7.6) {\iflangja{時間（秒）}{time (s)}};
\end{tikzpicture}
   (width ~116 mm)
\begin{tikzpicture}[x=1mm, y=1mm, font=\scriptsize\sffamily,
  >={Stealth[length=1.3mm]},
  row/.style={font=\scriptsize\sffamily, anchor=east},
  lab/.style={font=\scriptsize\sffamily, inner sep=0.8pt, align=center}]
  % time t (s) -> x = 18 + 1.4 t (mm)
  \node[row] at (16,16) {\iflangja{クライアント 1}{client 1}};
  \node[row] at (16,8) {\iflangja{サーバ}{server}};
  \node[row] at (16,0) {\iflangja{クライアント 2}{client 2}};
  \draw[ln-rule] (18,16) -- ({18+1.4*40},16);
  \draw[ln-rule, densely dotted] ({18+1.4*40},16) -- (116,16);
  \draw[ln-rule] (18,8) -- (116,8);
  \draw[ln-rule] (18,0) -- (116,0);
  % lock holder on the server
  \draw[fill=ln-box] ({18+1.4*0},6) rectangle ({18+1.4*55},10);
  \node[lab] at ({18+1.4*12},8) {\iflangja{クライアント 1 がロックを保持}{held by client 1}};
  \draw[fill=ln-accent!30] ({18+1.4*55},6) rectangle (116,10);
  \node[lab] at ({(18+1.4*55+116)/2},8) {\iflangja{クライアント 2}{client 2}};
  % lease expiry marks
  \foreach \t in {30,55} {
    \draw[densely dashed, ln-caution] ({18+1.4*\t},4.6) -- ({18+1.4*\t},11.4);
  }
  \node[lab, text=ln-caution, anchor=south] at ({18+1.4*55},11.4) {\iflangja{リース期限切れ}{lease expires}};
  \node[lab, text=ln-caution, anchor=north] at ({18+1.4*30},4.6) {\iflangja{当初の期限}{first expiry}};
  % client 1: acquire, renew, crash
  \draw[->, thick] ({18+1.4*0},14.8) -- ({18+1.4*0},10.3);
  \node[lab, anchor=south west] at ({18+1.4*0},17) {\iflangja{ロック獲得（リース 30 秒）}{lock (30 s lease)}};
  \draw[->, thick] ({18+1.4*25},14.8) -- ({18+1.4*25},10.3);
  \node[lab, anchor=south] at ({18+1.4*25},17) {\iflangja{更新}{renew}};
  \node[text=ln-caution, font=\normalsize\bfseries] at ({18+1.4*40},16) {$\times$};
  \node[lab, anchor=south, text=ln-caution] at ({18+1.4*40},17) {\iflangja{クラッシュ}{crash}};
  % client 2: request and grant
  \draw[->, thick] ({18+1.4*45},1.2) -- ({18+1.4*45},5.7);
  \node[lab, anchor=north] at ({18+1.4*45},-1.2) {\iflangja{ロック要求}{request}};
  \draw[densely dotted, thick] ({18+1.4*45},0) -- ({18+1.4*55},0);
  \draw[->, thick] ({18+1.4*55},5.7) -- ({18+1.4*55},1.2);
  \node[lab, anchor=north] at ({18+1.4*55},-1.2) {\iflangja{許可}{granted}};
  % time axis
  \draw[->] (18,-8) -- (117,-8);
  \foreach \t in {0,10,...,70} {
    \draw ({18+1.4*\t},-7.2) -- ({18+1.4*\t},-8.8);
    \node[lab, anchor=north] at ({18+1.4*\t},-9) {\t};
  }
  \node[lab, anchor=south east] at (117,-7.6) {\iflangja{時間（秒）}{time (s)}};
\end{tikzpicture}
   (width ~116 mm)
\begin{tikzpicture}[x=1mm, y=1mm, font=\scriptsize\sffamily,
  >={Stealth[length=1.3mm]},
  row/.style={font=\scriptsize\sffamily, anchor=east},
  lab/.style={font=\scriptsize\sffamily, inner sep=0.8pt, align=center}]
  % time t (s) -> x = 18 + 1.4 t (mm)
  \node[row] at (16,16) {\iflangja{クライアント 1}{client 1}};
  \node[row] at (16,8) {\iflangja{サーバ}{server}};
  \node[row] at (16,0) {\iflangja{クライアント 2}{client 2}};
  \draw[ln-rule] (18,16) -- ({18+1.4*40},16);
  \draw[ln-rule, densely dotted] ({18+1.4*40},16) -- (116,16);
  \draw[ln-rule] (18,8) -- (116,8);
  \draw[ln-rule] (18,0) -- (116,0);
  % lock holder on the server
  \draw[fill=ln-box] ({18+1.4*0},6) rectangle ({18+1.4*55},10);
  \node[lab] at ({18+1.4*12},8) {\iflangja{クライアント 1 がロックを保持}{held by client 1}};
  \draw[fill=ln-accent!30] ({18+1.4*55},6) rectangle (116,10);
  \node[lab] at ({(18+1.4*55+116)/2},8) {\iflangja{クライアント 2}{client 2}};
  % lease expiry marks
  \foreach \t in {30,55} {
    \draw[densely dashed, ln-caution] ({18+1.4*\t},4.6) -- ({18+1.4*\t},11.4);
  }
  \node[lab, text=ln-caution, anchor=south] at ({18+1.4*55},11.4) {\iflangja{リース期限切れ}{lease expires}};
  \node[lab, text=ln-caution, anchor=north] at ({18+1.4*30},4.6) {\iflangja{当初の期限}{first expiry}};
  % client 1: acquire, renew, crash
  \draw[->, thick] ({18+1.4*0},14.8) -- ({18+1.4*0},10.3);
  \node[lab, anchor=south west] at ({18+1.4*0},17) {\iflangja{ロック獲得（リース 30 秒）}{lock (30 s lease)}};
  \draw[->, thick] ({18+1.4*25},14.8) -- ({18+1.4*25},10.3);
  \node[lab, anchor=south] at ({18+1.4*25},17) {\iflangja{更新}{renew}};
  \node[text=ln-caution, font=\normalsize\bfseries] at ({18+1.4*40},16) {$\times$};
  \node[lab, anchor=south, text=ln-caution] at ({18+1.4*40},17) {\iflangja{クラッシュ}{crash}};
  % client 2: request and grant
  \draw[->, thick] ({18+1.4*45},1.2) -- ({18+1.4*45},5.7);
  \node[lab, anchor=north] at ({18+1.4*45},-1.2) {\iflangja{ロック要求}{request}};
  \draw[densely dotted, thick] ({18+1.4*45},0) -- ({18+1.4*55},0);
  \draw[->, thick] ({18+1.4*55},5.7) -- ({18+1.4*55},1.2);
  \node[lab, anchor=north] at ({18+1.4*55},-1.2) {\iflangja{許可}{granted}};
  % time axis
  \draw[->] (18,-8) -- (117,-8);
  \foreach \t in {0,10,...,70} {
    \draw ({18+1.4*\t},-7.2) -- ({18+1.4*\t},-8.8);
    \node[lab, anchor=north] at ({18+1.4*\t},-9) {\t};
  }
  \node[lab, anchor=south east] at (117,-7.6) {\iflangja{時間（秒）}{time (s)}};
\end{tikzpicture}

  \caption{リースに基づくロック．クライアント 1 は \SI{25}{\second} に
    リースを更新し，\SI{40}{\second} にクラッシュする．リースが
    \SI{55}{\second} に期限切れになると，サーバはクライアント 2 に
    ロックを与える．}
  \label{fig:l14-lease}
\end{figure}

\begin{example}
  あるサーバが \SI{30}{\second} のリースを与える．\margin{期間はプロトコル
  が定めるものではない．1つのサーバが与えるリースはすべて同じ長さであり，
  クライアントはそれをサーバの \texttt{lease\_time} 属性から読む
  \cite{rfc7530}．}クライアント~1 は
  $t=0$ にロックを獲得し，\SI{25}{\second} にリースを更新して
  \SI{55}{\second} まで延長するが，\SI{40}{\second} にクラッシュする．
  クライアント~2 は \SI{45}{\second} にロックを要求して待つ．
  \SI{55}{\second} にクライアント~1 のリースは更新されないまま期限切れと
  なり，サーバはクライアント~2 にロックを与える（\cref{fig:l14-lease}）．
  仮にクライアント~1 がクラッシュしたのではなく，単にサーバとの連絡を失った
  だけだったとしても，\SI{55}{\second} にはロックの使用をやめなければ
  ならない．ロックが別のクライアントに与えられたかどうかを知りえないから
  である．\caution{リースに時刻の同期は要らない．クライアントが計時を
  始めるのは付与が届いた後なので，常にサーバより先に権利を手放す．}
\end{example}

\begin{keypoint}
  NFS のクライアントは VFS，NFS クライアント，RPC を通じてリモートファイル
  に到達し，サーバが生成するファイルハンドルでファイルを識別する．NFSv3 は
  ステートレス，NFSv4 はステートフルである．NFS はセッションに基づく
  キャッシュと周期的な確認---既定でファイルは \SI{3}{\second} ごと，
  ディレクトリは \SI{30}{\second} ごと---を組み合わせる．NFSv4 はこれに
  委譲，リースに基づくロック，共有予約を加える．
\end{keypoint}

\section{Sprite 分散ファイルシステム}
\label{sec:l14-sprite}

分散オペレーティングシステム \term{Sprite}[Sprite] は 1980 年代に
カリフォルニア大学バークレー校で開発された．そのファイルシステムは，
ネットワーク上のワークステーションがサーバに格納されたファイルを共有
できるようにするものである．\source{P4L2，Sprite 分散ファイルシステム，
Sprite DFS のアクセスパターンの分析，Sprite DFS の分析から設計へ；
Nelson ら \cite{nelson1988}．}Nelson，Welch，Ousterhout
\cite{nelson1988} は，そのファイルシステムがどのようにデータをキャッシュ
するかを述べている．この論文が何より示唆に富むのはその設計過程である．
設計上の各決定が，ファイルが実際にどうアクセスされるかの分析から導かれ，
それによって正当化されている．\cref{tab:l14-sprite-analysis} に観測と，
そこから導かれた結論を示す．

\begin{table}[htbp]
  \centering
  \caption{アクセスパターンと，それが Sprite の設計にもたらした帰結．}
  \label{tab:l14-sprite-analysis}
  \small
  \begin{tabular}{@{}>{\raggedright\arraybackslash}p{0.4\linewidth}>{\raggedright\arraybackslash}p{0.54\linewidth}@{}}
    \toprule
    観測 & 帰結 \\
    \midrule
    全ファイルアクセスの \SI{33}{\percent} が書き込みである &
      アクセスの大半を占める読み出しについてはキャッシュが引き合う．
      しかし UNIXセマンティクスが要求するように，すべての書き込みを直ちに
      伝播するのは高価すぎる． \\
    ファイルの \SI{75}{\percent} は \SI{0.5}{\second} 未満，
      \SI{90}{\percent} は \SI{10}{\second} 未満しか open されない &
      セッションセマンティクスでもなおオーバーヘッドが大きすぎる．これほど
      短いセッションでは，open と close のたびの更新がきわめて頻繁に
      サーバへ届くことになる． \\
    新しいデータの
      \SIrange[range-phrase={--},range-units=single]{20}{30}{\percent} は
      \SI{30}{\second} 以内に，\SI{50}{\percent} は \SI{5}{\minute} 以内に
      削除される &
      close の時点でデータを書き戻す必要はない．データの多くは，他のどの
      クライアントもそれを必要とする前に削除される． \\
    ファイルへの並行アクセスはまれである & 並行アクセスは最適化する必要は
      ないが，支えなければならない． \\
    \bottomrule
  \end{tabular}
\end{table}

これらの観測から Sprite の設計が導かれる．
\begin{itemize}
  \item クライアントはファイルのブロックをキャッシュし，変更されたブロックを
    30 秒の遅延をもってサーバへ書き戻す．詳細は後述する．\margin{Linux も
    ページキャッシュに対して同じことをする．フラッシャスレッドが
    \SI{5}{\second} ごとに起き，\SI{30}{\second} を超えてダーティなままの
    データを書き戻す（\texttt{dirty\_expire\_centisecs}）．}
  \item 別のクライアントが書き込んでいたファイルをあるクライアントが
    open すると，サーバはまずそのクライアントからダーティなブロックを
    回収する．
  \item すべての open はサーバへ行き，ディレクトリはクライアントに
    キャッシュされない．
  \item ファイルが並行に書き込まれているときは，そのファイルのキャッシュを
    無効にし，そのファイルに対する操作はすべてサーバへ行く．
\end{itemize}

最初の項目が書き込み方針を定める．変更されたブロックをクライアントの
キャッシュに保持し，書き込みのたびでも close の時点でもなく，遅延の後に
はじめて書き戻すことを\term[ちえんかきこみ]{遅延書き込み}[delayed write]，
あるいは遅延書き戻しと呼ぶ．Sprite では，クライアントは 30 秒ごとに，
直前の 30 秒間に変更されなかったダーティなブロックを書き戻す．したがって，
まだ変更され続けているブロックが繰り返し転送されることはなく，書き込みから
30 秒以内に削除されたデータはまったく転送されない．ブロックがサーバに届く
のは，その最後の変更から 30 秒後から 60 秒後の間である．\caution{遅延書き
込みは通信量と引き換えに耐久性を手放す．クライアントがクラッシュすると，
まだ書き戻されていないブロックはすべて失われ，最大で1分ぶんの作業が失われ
る．Sprite がこれを受け入れるのは，そうしたブロックの多くがどのみち削除され
るからである．}

\begin{example}
  あるクライアントが $t = 0, 30, 60, 90, \ldots$~秒に書き戻すとする．
  ブロック $b_1$ は \SI{12}{\second} に変更される．\SI{30}{\second} の
  時点では最後の変更から \SI{18}{\second} しか経っていないので，書き戻され
  るのは \SI{60}{\second} の時点である．ブロック $b_2$ は \SI{12}{\second}
  と，さらに \SI{50}{\second} に変更される．\SI{60}{\second} の時点では
  最後の変更から \SI{10}{\second} しか経っていないので，書き戻されるのは
  \SI{90}{\second} の時点である．ブロック $b_3$ は，\SI{20}{\second} に
  書かれ \SI{45}{\second} に削除される一時ファイルに属しており，サーバへ
  送られることはない．
\end{example}

Sprite が提供するセマンティクスは，ファイルがどのように共有されるかに
依存する．

\begin{definition}[逐次書き込み共有と並行書き込み共有]
  あるクライアントがファイルを変更して close し，別のクライアントが後から
  それを open するとき，そのファイルは
  \term[ちくじかきこみきょうゆう]{逐次書き込み共有}[sequential write sharing]%
  の下にある．すなわち，複数のクライアントが次々にファイルを使う．ファイル
  が複数のクライアントで同時に open されており，そのうち少なくとも1つが
  書き込みのために open しているとき，そのファイルは
  \term[へいこうかきこみきょうゆう]{並行書き込み共有}[concurrent write sharing]%
  の下にある．
\end{definition}

逐次書き込み共有の下では Sprite はファイルをキャッシュし，それを open する
どのクライアントも，直前の書き手が書いたデータを見る．\caution{ここでの
\emph{逐次}はクライアントを指すのであって，ファイル内のアクセスの仕方を指す
のではない．ブロックを順に読むことではなく，クライアントが次々にファイルを
使うことである．}並行書き込み共有の
下では Sprite はファイルをキャッシュせず，サーバがすべての読み出しと
書き込みを順に実行する．これは UNIXセマンティクスを与える．

\section{Sprite におけるファイルアクセス操作}
\label{sec:l14-sprite-ops}

複数のクライアントが読み，1つ以上のクライアントが書くファイルを考える．
\source{P4L2，Sprite におけるファイルアクセス操作；Nelson ら
\cite{nelson1988}．}すべての open はサーバを経由し，ファイルがキャッシュ
可能である限りすべてのクライアントがブロックをキャッシュし，書き手は，
遅延書き込みが要求するとおり，変更した各ブロックについてタイムスタンプを
保持する．クライアントとサーバは，ファイルごとに次の状態を保持する
（\cref{fig:l14-sprite}）．
\begin{description}
  \item[クライアント] そのファイルについてキャッシュが有効かどうか，
    キャッシュしているブロック，ダーティな各ブロックの最終変更時刻，
    そしてキャッシュしているデータの版番号．
  \item[サーバ] そのファイルを読み出しのために open しているクライアントと
    書き込みのために open しているクライアント，最後にそのファイルに書き
    込んだクライアント，現在の版番号，そしてファイルがキャッシュ可能か
    どうか．
\end{description}

\begin{figure}[htbp]
  \centering
  % Sprite: per-file state on clients and server, and the messages when W3
% opens a file that W2 has open for writing (concurrent write sharing).
% Usage: % Sprite: per-file state on clients and server, and the messages when W3
% opens a file that W2 has open for writing (concurrent write sharing).
% Usage: % Sprite: per-file state on clients and server, and the messages when W3
% opens a file that W2 has open for writing (concurrent write sharing).
% Usage: \input{figures/l14-sprite}   (width ~118 mm)
\begin{tikzpicture}[x=1mm, y=1mm, font=\scriptsize\sffamily,
  >={Stealth[length=1.5mm]},
  machine/.style={draw, rounded corners=0.8mm},
  st/.style={font=\scriptsize\sffamily, anchor=north west, inner sep=0pt},
  lab/.style={font=\scriptsize\sffamily, align=center, inner sep=0.8pt},
  who/.style={font=\scriptsize\sffamily\bfseries, anchor=north west}]
  % client W2 (writer)
  \draw[machine, fill=white] (0,27) rectangle (46,47);
  \node[who] at (1,46.5) {\iflangja{クライアント $W_2$（書き込み中）}{client $W_2$ (writing)}};
  \node[st] at (2,41.5) {\begin{tabular}{@{}l@{\quad}l@{}}
      \iflangja{キャッシュ}{caching} & \iflangja{有効 → 無効}{on $\to$ off}\\
      \iflangja{ブロック}{blocks} & 0, 1, 5\\
      \iflangja{ダーティ}{dirty} & 1 (41\,s), 5 (44\,s)\\
      \iflangja{版}{version} & 6
    \end{tabular}};
  % client W3 (opening)
  \draw[machine, fill=white] (72,27) rectangle (118,47);
  \node[who] at (73,46.5) {\iflangja{クライアント $W_3$（open 中）}{client $W_3$ (opening)}};
  \node[st] at (74,41.5) {\begin{tabular}{@{}l@{\quad}l@{}}
      \iflangja{キャッシュ}{caching} & \iflangja{無効}{off}\\
      \iflangja{ブロック}{blocks} & ---\\
      \iflangja{版}{version} & 7
    \end{tabular}};
  % server
  \draw[machine, fill=ln-box] (33,-8) rectangle (85,13);
  \node[who] at (34,12.5) {\iflangja{サーバ}{server}};
  \node[st] at (35,7.5) {\begin{tabular}{@{}l@{\quad}l@{}}
      \iflangja{読み手}{readers} & ---\\
      \iflangja{書き手}{writers} & $W_2$, $W_3$\\
      \iflangja{版}{version} & $6 \to 7$\\
      \iflangja{キャッシュ可能}{cacheable} & \iflangja{可 → 不可}{yes $\to$ no}
    \end{tabular}};
  % (1) open by W3
  \draw[->, thick] (108,27) |- (85,2);
  \node[lab, anchor=south] at (96,2.6) {\iflangja{(1) open\\（書き込み）}{(1) open\\for writing}};
  % (2) dirty blocks from W2
  \draw[<->, thick, ln-accent] (33,2) -| (12,27);
  \node[lab, anchor=west, text=ln-accent] at (13,16.5) {\iflangja{(2) ダーティ\\ブロック}{(2) dirty\\blocks}};
  % (3) disable caching
  \draw[->, densely dashed, ln-caution, thick] (59,13) -- (59,20) -| (36,27);
  \draw[->, densely dashed, ln-caution, thick] (59,20) -| (82,27);
  \node[lab, anchor=south, text=ln-caution, fill=white] at (59,20.6) {\iflangja{(3) キャッシュ無効化}{(3) disable caching}};
\end{tikzpicture}
   (width ~118 mm)
\begin{tikzpicture}[x=1mm, y=1mm, font=\scriptsize\sffamily,
  >={Stealth[length=1.5mm]},
  machine/.style={draw, rounded corners=0.8mm},
  st/.style={font=\scriptsize\sffamily, anchor=north west, inner sep=0pt},
  lab/.style={font=\scriptsize\sffamily, align=center, inner sep=0.8pt},
  who/.style={font=\scriptsize\sffamily\bfseries, anchor=north west}]
  % client W2 (writer)
  \draw[machine, fill=white] (0,27) rectangle (46,47);
  \node[who] at (1,46.5) {\iflangja{クライアント $W_2$（書き込み中）}{client $W_2$ (writing)}};
  \node[st] at (2,41.5) {\begin{tabular}{@{}l@{\quad}l@{}}
      \iflangja{キャッシュ}{caching} & \iflangja{有効 → 無効}{on $\to$ off}\\
      \iflangja{ブロック}{blocks} & 0, 1, 5\\
      \iflangja{ダーティ}{dirty} & 1 (41\,s), 5 (44\,s)\\
      \iflangja{版}{version} & 6
    \end{tabular}};
  % client W3 (opening)
  \draw[machine, fill=white] (72,27) rectangle (118,47);
  \node[who] at (73,46.5) {\iflangja{クライアント $W_3$（open 中）}{client $W_3$ (opening)}};
  \node[st] at (74,41.5) {\begin{tabular}{@{}l@{\quad}l@{}}
      \iflangja{キャッシュ}{caching} & \iflangja{無効}{off}\\
      \iflangja{ブロック}{blocks} & ---\\
      \iflangja{版}{version} & 7
    \end{tabular}};
  % server
  \draw[machine, fill=ln-box] (33,-8) rectangle (85,13);
  \node[who] at (34,12.5) {\iflangja{サーバ}{server}};
  \node[st] at (35,7.5) {\begin{tabular}{@{}l@{\quad}l@{}}
      \iflangja{読み手}{readers} & ---\\
      \iflangja{書き手}{writers} & $W_2$, $W_3$\\
      \iflangja{版}{version} & $6 \to 7$\\
      \iflangja{キャッシュ可能}{cacheable} & \iflangja{可 → 不可}{yes $\to$ no}
    \end{tabular}};
  % (1) open by W3
  \draw[->, thick] (108,27) |- (85,2);
  \node[lab, anchor=south] at (96,2.6) {\iflangja{(1) open\\（書き込み）}{(1) open\\for writing}};
  % (2) dirty blocks from W2
  \draw[<->, thick, ln-accent] (33,2) -| (12,27);
  \node[lab, anchor=west, text=ln-accent] at (13,16.5) {\iflangja{(2) ダーティ\\ブロック}{(2) dirty\\blocks}};
  % (3) disable caching
  \draw[->, densely dashed, ln-caution, thick] (59,13) -- (59,20) -| (36,27);
  \draw[->, densely dashed, ln-caution, thick] (59,20) -| (82,27);
  \node[lab, anchor=south, text=ln-caution, fill=white] at (59,20.6) {\iflangja{(3) キャッシュ無効化}{(3) disable caching}};
\end{tikzpicture}
   (width ~118 mm)
\begin{tikzpicture}[x=1mm, y=1mm, font=\scriptsize\sffamily,
  >={Stealth[length=1.5mm]},
  machine/.style={draw, rounded corners=0.8mm},
  st/.style={font=\scriptsize\sffamily, anchor=north west, inner sep=0pt},
  lab/.style={font=\scriptsize\sffamily, align=center, inner sep=0.8pt},
  who/.style={font=\scriptsize\sffamily\bfseries, anchor=north west}]
  % client W2 (writer)
  \draw[machine, fill=white] (0,27) rectangle (46,47);
  \node[who] at (1,46.5) {\iflangja{クライアント $W_2$（書き込み中）}{client $W_2$ (writing)}};
  \node[st] at (2,41.5) {\begin{tabular}{@{}l@{\quad}l@{}}
      \iflangja{キャッシュ}{caching} & \iflangja{有効 → 無効}{on $\to$ off}\\
      \iflangja{ブロック}{blocks} & 0, 1, 5\\
      \iflangja{ダーティ}{dirty} & 1 (41\,s), 5 (44\,s)\\
      \iflangja{版}{version} & 6
    \end{tabular}};
  % client W3 (opening)
  \draw[machine, fill=white] (72,27) rectangle (118,47);
  \node[who] at (73,46.5) {\iflangja{クライアント $W_3$（open 中）}{client $W_3$ (opening)}};
  \node[st] at (74,41.5) {\begin{tabular}{@{}l@{\quad}l@{}}
      \iflangja{キャッシュ}{caching} & \iflangja{無効}{off}\\
      \iflangja{ブロック}{blocks} & ---\\
      \iflangja{版}{version} & 7
    \end{tabular}};
  % server
  \draw[machine, fill=ln-box] (33,-8) rectangle (85,13);
  \node[who] at (34,12.5) {\iflangja{サーバ}{server}};
  \node[st] at (35,7.5) {\begin{tabular}{@{}l@{\quad}l@{}}
      \iflangja{読み手}{readers} & ---\\
      \iflangja{書き手}{writers} & $W_2$, $W_3$\\
      \iflangja{版}{version} & $6 \to 7$\\
      \iflangja{キャッシュ可能}{cacheable} & \iflangja{可 → 不可}{yes $\to$ no}
    \end{tabular}};
  % (1) open by W3
  \draw[->, thick] (108,27) |- (85,2);
  \node[lab, anchor=south] at (96,2.6) {\iflangja{(1) open\\（書き込み）}{(1) open\\for writing}};
  % (2) dirty blocks from W2
  \draw[<->, thick, ln-accent] (33,2) -| (12,27);
  \node[lab, anchor=west, text=ln-accent] at (13,16.5) {\iflangja{(2) ダーティ\\ブロック}{(2) dirty\\blocks}};
  % (3) disable caching
  \draw[->, densely dashed, ln-caution, thick] (59,13) -- (59,20) -| (36,27);
  \draw[->, densely dashed, ln-caution, thick] (59,20) -| (82,27);
  \node[lab, anchor=south, text=ln-caution, fill=white] at (59,20.6) {\iflangja{(3) キャッシュ無効化}{(3) disable caching}};
\end{tikzpicture}

  \caption{Sprite のファイルごとの状態．$W_2$ が書き込みのために open して
    いるファイルを $W_3$ が open すると (1)，サーバは $W_2$ のダーティな
    ブロックを回収し (2)，両クライアントのキャッシュを無効にする (3)．}
  \label{fig:l14-sprite}
\end{figure}

\paragraph{逐次書き込み共有．}クライアントがファイルを書き込みのために
open するたびに，サーバは版番号を増やす．ファイルを open したクライアントは，
サーバが返した版番号を自分のキャッシュしたブロックの版番号と比較し，古け
ればそのブロックを捨てる．遅延書き込みのため，最後にファイルに書き込んだ
クライアントは，close した後もなおそのファイルのダーティなブロックを
保持していることがある．そこでサーバは，別のクライアントがファイルを
open する前にそのクライアントにそれらを書き戻すよう求め，新しいクライアント
が最新のデータを見られるようにする．キャッシュは有効なままである．

\paragraph{並行書き込み共有．}別のクライアントが open しているファイルを
あるクライアントが書き込みのために open したとき，あるいは別のクライアント
が書き込みのために open しているファイルを読み出しのために open したとき，
サーバは再び書き手のダーティなブロックを回収し，ファイルをキャッシュ不可と
印し，それを open しているすべてのクライアントに通知する．それ以降，どの
クライアントもそのファイルをキャッシュせず，すべての読み出しと書き込みが
サーバへ行く．すべてのクライアントがファイルを close すると，ファイルは
再びキャッシュ可能になる．

\begin{example}
  \cref{tab:l14-sprite-trace} は，はじめ版~4 にあるファイル $F$ について，
  クライアント $R_1$，$W_1$，$W_2$，$W_3$ が open と close を行う間の
  サーバの状態を追ったものである．段階~5 では $W_1$ は $F$ を close して
  いるがなおダーティなブロックを保持しており，サーバは $W_2$ を進める前に
  それを回収する．共有は逐次的であり，$W_2$ は $F$ をキャッシュしてよい．
  \tip{逐次共有でキャッシュが生き残るのは，サーバが先にブロックを回収できる
  「最後の書き手」が必ずいるからである．並行共有にはそれがいないので，サーバ
  がすべてのアクセスを自ら処理するほかない．}
  段階~6 では $W_2$ がまだ書き込みのために open している $F$ を $W_3$ が
  open し（\cref{fig:l14-sprite}），キャッシュが無効にされる．段階~7 では
  $R_1$ が，キャッシュしている版~4 のブロックが古いことを見出し，それを
  捨ててサーバから読む．
\end{example}

\begin{table}[htbp]
  \centering
  \caption{Sprite におけるファイル $F$ のサーバ状態．版とキャッシュ可否は
    各段階の後の値を示す．}
  \label{tab:l14-sprite-trace}
  \small
  \begin{tabular}{@{}r>{\raggedright\arraybackslash}p{0.24\linewidth}>{\raggedright\arraybackslash}p{0.29\linewidth}llc@{}}
    \toprule
    & 事象 & サーバの動作 & open 中 & 版 & キャッシュ \\
    \midrule
    1 & $R_1$ が読み出しで open & --- & $R_1$ & 4 & 可 \\
    2 & $R_1$ が close & --- & --- & 4 & 可 \\
    3 & $W_1$ が書き込みで open & 版を増やす & $W_1$ & 5 & 可 \\
    4 & $W_1$ が書き込み，close & ---（ダーティなブロックは $W_1$ に残る）
      & --- & 5 & 可 \\
    5 & $W_2$ が書き込みで open & $W_1$ のダーティなブロックを回収；
      版を増やす & $W_2$ & 6 & 可 \\
    6 & $W_3$ が書き込みで open & $W_2$ のダーティなブロックを回収；
      版を増やす；キャッシュを無効化 & $W_2$, $W_3$ & 7 & 不可 \\
    7 & $R_1$ が読み出しで open & $R_1$ にキャッシュしないよう伝える &
      $R_1$, $W_2$, $W_3$ & 7 & 不可 \\
    8 & $W_2$, $W_3$, $R_1$ が close & キャッシュを有効化 & --- & 7 & 可 \\
    \bottomrule
  \end{tabular}
\end{table}

\needspace{18\baselineskip}
\begin{keypoint}[章のまとめ]
  分散ファイルシステムは，ネットワークで結ばれたマシン上のファイルを，
  通常のファイルシステムインタフェースを通じて提供する．その構成は，単一の
  サーバ，複製あるいは分割されたサーバ，そしてピアのいずれかである．実用的な
  設計はアップロード・ダウンロードモデルと遠隔アクセスモデルの間に位置する．
  クライアントはブロックをキャッシュし，ステートフルサーバと頻繁にやり取り
  し，サーバはキャッシュされた複製の整合を保たなければならない．ファイル
  共有セマンティクス---UNIXセマンティクス，セッションセマンティクス，
  リースを伴う周期的な更新，不変ファイル，トランザクション---がクライアント
  の観測するものを定め，ファイルとディレクトリで異なってよい．レプリケー
  ションは書き込みが複雑になる代価と引き換えに可用性と負荷の均衡を高め，
  パーティショニングはファイルシステムの規模に対してスケールする．NFS は
  RPC の上でファイルハンドルを用い，ステートレスな NFSv3 かステートフルな
  NFSv4 で，セッションに基づくキャッシュと周期的な確認，リースに基づく
  ロックを提供する．Sprite は測定されたアクセスパターンから設計を導いた．
  すなわち，30 秒の遅延を伴う遅延書き込み，逐次書き込み共有の下での
  キャッシュ，並行書き込み共有の下でのキャッシュの停止である．
\end{keypoint}

\end{document}
